\documentclass[11pt]{amsart}

\usepackage[T1]{fontenc}
\usepackage[utf8]{inputenc}
\usepackage{lmodern}
\usepackage{microtype}
\usepackage[margin=1.15in]{geometry}
\usepackage{amsmath,amssymb,mathtools,amscd}
\usepackage{mathrsfs}
\usepackage{amsthm}
\usepackage{booktabs}
\usepackage{enumitem}
\usepackage{xcolor}
\usepackage[colorlinks=true,linkcolor=blue!55!black,citecolor=green!45!black,urlcolor=blue!55!black]{hyperref}

\newtheorem{theorem}{Theorem}[section]
\newtheorem{proposition}[theorem]{Proposition}
\newtheorem{lemma}[theorem]{Lemma}
\newtheorem{corollary}[theorem]{Corollary}
\newtheorem{definition}[theorem]{Definition}
\newtheorem{example}[theorem]{Example}
\newtheorem{remark}[theorem]{Remark}
\newtheorem{problem}[theorem]{Problem}

\newcommand{\End}{\operatorname{End}}
\newcommand{\Hom}{\operatorname{Hom}}
\newcommand{\Span}{\operatorname{span}}
\newcommand{\Spec}{\operatorname{Spec}}
\newcommand{\coker}{\operatorname{coker}}
\newcommand{\rank}{\operatorname{rank}}
\newcommand{\Obs}{\operatorname{Obs}}
\newcommand{\Wfib}{\mathfrak W}
\newcommand{\Rfib}{\mathfrak R}
\newcommand{\cO}{\mathcal O}
\newcommand{\cD}{\mathscr D}
\newcommand{\cQ}{\mathscr Q}
\newcommand{\cP}{\mathscr P}
\newcommand{\C}{\mathbb C}
\newcommand{\Q}{\mathbb Q}
\newcommand{\Spf}{\operatorname{Spf}}
\newcommand{\MC}{\operatorname{MC}}
\newcommand{\gpres}{\mathfrak p}
\newcommand{\ghoch}{\mathfrak g}
\newcommand{\Rel}{\mathscr C}

\title[The Inverse Leibniz Problem]{The Inverse Leibniz Problem:\\Reconstruction Fibers, Rigidity, Obstructions, and Deformation DGLAs}
\author{Akari H.}
\date{Progress-preservation draft v0.41 --- August 2026}

\begin{document}

\begin{abstract}
Let $A$ be an associative algebra, $V$ a vector space, and $B:A\times A\to V$ a bilinear map.  A Leibniz identity for $B$ is typed only after left and right actions of $A$ on $V$ have been specified.  We study the inverse problem of reconstructing those actions from $B$.  Weak reconstruction is affine-linear and is governed by a Leibniz observability operator, whereas strict reconstruction intersects the weak affine fiber with the quadratic equations for a unital bimodule.  Finite-dimensional examples show that spanning of $V$ by $B(A,A)$ does not imply global identifiability: one example has four globally distinct but infinitesimally and locally rigid strict reconstructions.  Linearization yields a strict tangent observability operator and a second-order obstruction map, and a dual-number example exhibits a nonzero first-order direction killed at second order.

The translated strict equations admit both a finite presentation DGLA and an intrinsic constrained Hochschild DGLA.  In a cubic example with $A=\C[\epsilon]/(\epsilon^4)$ and $V=\C^3$, the strict tangent space is two-dimensional, the quadratic Kuranishi map vanishes, and a cubic obstruction survives.  Exact elimination gives the complete formal germ
\[
\widehat{\cO}_{\mathbf R_B,p}\cong \C[[u,v]]/(v^3).
\]
A presentation-level quartic transfer obstruction is shown to be truncation-sensitive: it cannot be removed by changing a contraction of the fixed presentation DGLA, but it disappears in the intrinsic constrained Hochschild model.  Exact intrinsic homotopy transfer is then continued through higher arities.

The higher-order computation develops into a finite-feedback problem.  Pure-mixed degree-one cochains provide a split contractible reservoir; finite controller boxes remain homologically healthy while their actuator reachability undergoes bifurcations.  In the support block $(\epsilon^2,\epsilon)$, the quotient maps $T_\alpha(u)=d[\alpha,u]$ and $T_\xi(u)=d[\xi,u]$ determine an exact actuator flag with inert, one-sided, and genuinely two-sided strata.  This flag correctly predicts two successive fresh actuators.  From arity $19$ onward the dynamics changes phase: the same finite box $U_{16}$ controls arities $19,20,21,22$ without enlargement by reusing $E_{32}$, while $E_{33}$ remains a reserve two-sided actuator.

Finally, a $v$-degree filtration shows that the stationary $E/K/P$ feedback variables are terminal corrections: they cannot carry information into future distinguished $v^6$ forcing.  The true memory is an autonomous low-rail subsystem $x_b(k)=F_{k+b}^{(k,b)}$, $b=2,3,4,5$, satisfying exact triangular quadratic recursions.  A frozen order-$14$ transfer law passes two genuine out-of-sample tests at arities $22$ and $23$.  More importantly, the homogeneous operator $A=-h\,\mathrm{ad}_\alpha$ and the triangular Cauchy forcing admit an all-orders spectral analysis.  With
\[
E=S+2,
\qquad
q(S)=S^5-4S^4+12S^3-32S^2+80S-192,
\]
the stabilized homogeneous module has minimal polynomial $m_A=Eq$, while the low rails have
\[
p_3=E^2q,
\qquad
p_4=E^3q,
\qquad
p_5=E^4q^2.
\]
Residual modules, finite landing operators, and gcd formulas then give a complete spectral stratification of the $x_5$ extension.  The homogeneous landing, forcing-memory, and canonical higher-lift layers are classified respectively by $\gcd(m_A,r_d)$, $\gcd(p_4,r_d)$, and $\gcd(p_5,r_d)$.  Over a splitting field the finest canonical stratification has $708$ nonzero strata plus the zero datum; over $\Q$ it collapses to eight nonzero strata plus zero.  Gauge quotients, general triangular spectral-extension theorems beyond the fixed cubic contraction, and definitive literature positioning remain open.
\end{abstract}

\maketitle
\tableofcontents

\section{Introduction}

Classical deformation theory typically starts with an algebraic structure and studies its infinitesimal and formal variations.  In the present problem, the direction is partly reversed.  We are given a bilinear map
\[
B:A\times A\longrightarrow V
\]
and ask whether the left and right $A$-actions on $V$ can be recovered from the requirement that $B$ satisfy Leibniz rules in both variables.

The key typing observation is elementary but decisive: a formula such as
\[
B(ab,c)=a\,B(b,c)+B(a,c)\,b
\]
is meaningless unless the products between $A$ and $V$ have already been defined.  Thus ``$B$ is a biderivation'' is a property of a triple $(B,L,R)$, not of $B$ alone.  The inverse Leibniz problem asks how much of $(L,R)$ is observable from $B$.

The present draft establishes twelve layers of the problem.
\begin{enumerate}[label=(\roman*)]
\item Weak reconstruction is linear and is controlled by a Leibniz observability operator.
\item Strict reconstruction imposes the unital bimodule equations and is generally nonlinear.
\item Linearization at a strict reconstruction separates infinitesimal, local, and global rigidity.
\item A quadratic obstruction class decides whether a strict tangent vector lifts to second order.
\item The full translated equations admit both a finite presentation DGLA and an intrinsic constrained Hochschild DGLA formulation.
\item The Kuranishi ideal determines an intrinsic minimal relation space whose dual measures the obstruction directions that actually generate the local equations.
\item Higher Kuranishi operations can occur intrinsically: an explicit example has vanishing quadratic Kuranishi map but nonzero cubic term.
\item A strict contraction need not realize the simplest formal Kuranishi normal form in the finite presentation DGLA: every presentation transfer in the cubic example has a nonzero quartic operation even though the formal relation ideal is equivalent to the pure cubic equation $v^3=0$.
\item This quartic no-go is not intrinsic to the constrained Hochschild DGLA.  We construct an intrinsic contraction with the same cubic obstruction but with $\ell_4|_{\operatorname{Sym}^4H^1}=0$, showing that the degree-three Hochschild differential removes a presentation-level normal-form obstruction.
\item The same intrinsic contraction can be extended through arity five: explicit quartic homotopies and a rank-five chain-level certificate give $\ell_5|_{\operatorname{Sym}^5H^1}=0$ without changing the nonzero cubic obstruction.
\item The contraction extends one step further through arity six: three explicit quintic homotopy lifts and a rank-nine chain-level certificate give $\ell_6|_{\operatorname{Sym}^6H^1}=0$ simultaneously with the quartic and quintic vanishings.
\item The same contraction extends through arity seven: four pure-mixed sextic homotopy lifts enlarge the degree-one splitting to rank eleven, and a rank-fourteen chain-level certificate gives $\ell_7|_{\operatorname{Sym}^7H^1}=0$ simultaneously with all previous higher vanishings.
\item A split pure-mixed reservoir and finite-feedback analysis continue the same intrinsic contraction beyond arity seven.  Finite controller boxes remain acyclic but undergo exact reachability bifurcations, beginning at arity thirteen.
\item The $(\epsilon^2,\epsilon)$ mixed block admits a block actuator flag theorem.  Its two-sided stratum evolves by deletion under successive rescued histories, and this flag predicts the fresh actuators $E_{31}$ and $E_{32}$ before the corresponding arity computations.
\item At arity nineteen the dynamics changes from fresh-actuator growth to stationary reuse.  One fixed box $U_{16}$ controls arities nineteen through twenty-two, with a stationary $E_{32}$ actuator sensitivity and no use of the reserved $E_{33}$ direction.
\item A $v$-degree filtration proves controller terminality.  The future forcing memory lies instead in four autonomous low-$v$ rails.  A frozen order-fourteen transfer realization passes two out-of-sample tests, after which the same recurrence is proved all-orders from the homogeneous operator $A=-h\,\mathrm{ad}_\alpha$ and triangular Cauchy forcing.
\item The residual modules satisfy
\[
\mathcal R_3\simeq J_{-2}^{(1)},
\qquad
\mathcal R_4\simeq J_{-2}^{(2)},
\qquad
\mu_{\mathcal R_5}=E^3q.
\]
This yields a Triangular Spectral Extension Lemma that separates forcing annihilation from the homogeneous spectral module reached by the residual.
\item Complete finite landing operators determine the next spectral multiplier without constructing an entire temporal rail.  For $x_4$ the landing image is contained in $\ker(A+2I)$, forcing an $E$-extension; for $x_5$ the landing operator is surjective onto the six-dimensional homogeneous module and full $m_A$-extension is Zariski-generic.
\item The exceptional landing geometry admits exact gcd stratifications.  The three nested invariants
\[
\gcd(m_A,r_d),\qquad \gcd(p_4,r_d),\qquad \gcd(p_5,r_d)
\]
classify homogeneous spectral support, forcing-quotient memory, and canonical higher-lift contact.  Over a splitting field the finest nonzero strata are indexed by truncated spectral multijets of total contact weight at most seven.
\end{enumerate}

The framework is adjacent to, but directionally different from, universal differential calculi and double-derivation formalisms, where a target bimodule structure is normally fixed in advance.  It also follows the standard deformation-theoretic pattern in which tangent vectors are cocycle-like data and higher-order extension is controlled by obstruction classes \cite{Gerstenhaber1964,Schlessinger1968}.  A full comparison with universal noncommutative differentials and double Poisson geometry \cite{CuntzQuillen1995,VanDenBergh2008} is deferred.

\subsection*{Status and limitations}
This is a progress-preservation draft rather than a submission-ready manuscript.  The principal finite-dimensional reconstruction, obstruction, DGLA, finite-feedback, block-actuator, stationary-reuse, and low-rail calculations have been checked by exact integer or rational symbolic computation, with modular minors used only as lower-bound certificates when paired with exact rational upper bounds.  The cubic example is verified through its complete formal germ, through full harmonic feedback at arity twenty-two, through the autonomous low-rail recursion at arity twenty-three, and through the all-orders autonomous low-rail theorem for the fixed contraction.  The order-fourteen transfer law has two genuine frozen out-of-sample validations and is subsequently derived from the stabilized homogeneous operator and triangular forcing, rather than inferred from those data.  The landing and extension stratifications are exact for the canonical polynomial representatives in the fixed cubic $x_5$ module.  No claim of literature-level novelty is made yet.  The main remaining tasks are a systematic literature review, treatment of gauge equivalence, abstraction of the triangular spectral-extension mechanism to general contraction data, and extension to simultaneous deformations of $A$ or $B$.

\section{Typed reconstruction data}

Throughout, $A$ is a unital associative algebra over $\C$ and $V$ is a complex vector space.  Fix a bilinear map
\[
B:A\times A\longrightarrow V.
\]
A candidate left-right action pair consists of linear maps
\[
L:A\longrightarrow\End(V),
\qquad
R:A\longrightarrow\End(V).
\]
We use the notation
\[
a\triangleright v=L(a)v,
\qquad
v\triangleleft a=R(a)v.
\]
No multiplicativity is assumed at first.

\begin{definition}[Weak reconstruction fiber]
The weak reconstruction fiber $\Wfib(B)$ is the set of pairs $(L,R)$ satisfying
\begin{align}
B(ab,c)&=L(a)B(b,c)+R(b)B(a,c), \label{eq:weak-left}\\
B(a,bc)&=R(c)B(a,b)+L(b)B(a,c) \label{eq:weak-right}
\end{align}
for all $a,b,c\in A$.
\end{definition}

\begin{definition}[Strict reconstruction fiber]
The strict reconstruction fiber $\Rfib(B)$ is the subset of $\Wfib(B)$ consisting of pairs satisfying
\begin{align}
L(1)&=R(1)=I_V, \label{eq:unit}\\
L(ab)&=L(a)L(b), \label{eq:left-rep}\\
R(ab)&=R(b)R(a), \label{eq:right-rep}\\
[L(a),R(b)]&=0 \label{eq:commuting-actions}
\end{align}
for all $a,b\in A$.
\end{definition}

Thus a point of $\Rfib(B)$ makes $V$ into a unital $A$-bimodule for which $B$ is a biderivation relative to the convention in \eqref{eq:weak-left}--\eqref{eq:weak-right}.

\section{The Leibniz observability operator}

Suppose $(L,R)$ and $(L',R')$ lie in $\Wfib(B)$, and set
\[
X=L'-L,
\qquad
Y=R'-R.
\]
Subtracting the two weak Leibniz systems gives
\begin{align}
X(a)B(b,c)+Y(b)B(a,c)&=0, \label{eq:obs1}\\
Y(c)B(a,b)+X(b)B(a,c)&=0. \label{eq:obs2}
\end{align}

\begin{definition}[Leibniz observability operator]
Define
\[
\cO_B:\Hom(A,\End V)^2\longrightarrow \Hom(A^{\otimes3},V)^2
\]
by
\[
\cO_B(X,Y)(a,b,c)=
\begin{pmatrix}
X(a)B(b,c)+Y(b)B(a,c)\\
Y(c)B(a,b)+X(b)B(a,c)
\end{pmatrix}.
\]
\end{definition}

\begin{proposition}[Affine structure of the weak fiber]\label{prop:weak-affine}
If $\Wfib(B)$ is nonempty and $(L_0,R_0)\in\Wfib(B)$, then
\[
\Wfib(B)=(L_0,R_0)+\ker\cO_B.
\]
In particular, weak reconstruction is unique if and only if $\ker\cO_B=0$.
\end{proposition}

\begin{proof}
Write $(L,R)=(L_0+X,R_0+Y)$.  Equations \eqref{eq:weak-left}--\eqref{eq:weak-right} are affine-linear in $(L,R)$, and their homogeneous part is exactly \eqref{eq:obs1}--\eqref{eq:obs2}.
\end{proof}

\section{Finite-dimensional matrix realization}

Assume $\dim A=m$ and $\dim V=d$.  Choose bases $e_1,\dots,e_m$ of $A$ and $v_1,\dots,v_d$ of $V$.  Write
\[
e_i e_j=\sum_{\ell=1}^m C_{ij}^{\ell}e_\ell,
\qquad
B(e_i,e_j)=\sum_{\beta=1}^d B_{ij}^{\beta}v_\beta.
\]
Let $L_i=L(e_i)$ and $R_i=R(e_i)$.  The weak equations become
\begin{align}
\sum_{\beta}(L_i)_{\alpha\beta}B_{jk}^{\beta}
+\sum_{\beta}(R_j)_{\alpha\beta}B_{ik}^{\beta}
&=\sum_{\ell}C_{ij}^{\ell}B_{\ell k}^{\alpha}, \label{eq:matrix-left}\\
\sum_{\beta}(R_k)_{\alpha\beta}B_{ij}^{\beta}
+\sum_{\beta}(L_j)_{\alpha\beta}B_{ik}^{\beta}
&=\sum_{\ell}C_{jk}^{\ell}B_{i\ell}^{\alpha}. \label{eq:matrix-right}
\end{align}
After ordering all matrix entries, these equations give
\[
M_Bu=r_B,
\qquad
M_B\in M_{2m^3d\,\times\,2md^2}(\C).
\]
Consequently,
\[
\Wfib(B)\neq\varnothing
\iff
\rank M_B=\rank[M_B\mid r_B],
\]
and, if consistent,
\[
\dim\Wfib(B)=2md^2-\rank M_B.
\]

The strict equations add the quadratic constraints
\begin{align*}
L_iL_j&=\sum_{\ell}C_{ij}^{\ell}L_\ell,\\
R_jR_i&=\sum_{\ell}C_{ij}^{\ell}R_\ell,\\
L_iR_j&=R_jL_i,
\end{align*}
together with the unit equations.  Hence in finite dimension the strict reconstruction space is naturally an affine scheme cut out by a linear system and quadratic bimodule equations.

\section{A finite-dimensional non-identifiability theorem}\label{sec:finite-example}

We now show that the spanning condition $V=\Span B(A,A)$ does not force strict uniqueness.

Let
\[
q(x)=(x-1)(x-2)(x-3)(x-4)
\]
and
\[
A=\C[x]/(q(x)).
\]
Let $V=\C^2$, $v_0=(1,1)^T$, and define
\[
L_x=\begin{pmatrix}1&0\\0&2\end{pmatrix},
\qquad
R_x=\begin{pmatrix}3&0\\0&4\end{pmatrix}.
\]
For a polynomial class $p\in A$, set
\[
\Delta_{L,R}(p)=\bigl(p(L_x)-p(R_x)\bigr)(L_x-R_x)^{-1}
\]
and define
\[
B(p,q)=\Delta_{L,R}(p)\Delta_{L,R}(q)v_0.
\]
Since all matrices involved are diagonal, the divided-difference identity
\[
\Delta_{L,R}(pq)=p(L_x)\Delta_{L,R}(q)+\Delta_{L,R}(p)q(R_x)
\]
implies both weak Leibniz identities.

Moreover,
\[
B(x,x)=\begin{pmatrix}1\\1\end{pmatrix},
\qquad
B(x^2,x)=\begin{pmatrix}4\\6\end{pmatrix},
\]
which are linearly independent.  Thus
\[
V=\Span B(A,A).
\]

The scalar divided difference is symmetric in its two evaluation points.  Consequently, independently swapping the two points in either coordinate of $V$ leaves $B$ unchanged.  This produces the four strict reconstructions
\[
\begin{array}{c|c}
L_x&R_x\\ \hline
\operatorname{diag}(1,2)&\operatorname{diag}(3,4)\\
\operatorname{diag}(3,2)&\operatorname{diag}(1,4)\\
\operatorname{diag}(1,4)&\operatorname{diag}(3,2)\\
\operatorname{diag}(3,4)&\operatorname{diag}(1,2).
\end{array}
\]

\begin{theorem}[Global non-identifiability]\label{thm:four-points}
There exist finite-dimensional $A,V$, and $B:A\times A\to V$ such that
\[
V=\Span B(A,A)
\]
but $|\Rfib(B)|=4$.  Hence the spanning condition alone does not imply unique reconstruction.
\end{theorem}

\begin{proof}
The construction above gives at least four strict reconstructions.  Exact elimination in the basis $1,x,x^2,x^3$ gives a weak system with $32$ unknowns and rank $28$.  Parameterizing the weak solution by the entries of $R(x^3)$, the strict equations have Groebner basis
\[
(R(x^3)_{11}-1)(R(x^3)_{11}-27),\quad
R(x^3)_{12},\quad R(x^3)_{21},\quad
(R(x^3)_{22}-8)(R(x^3)_{22}-64).
\]
Thus there are exactly four strict points, corresponding to the table above.
\end{proof}

\section{Strict tangent observability and rigidity}

Fix $p=(L,R)\in\Rfib(B)$.  Define
\begin{align*}
d_LX(a,b)&=X(a)L(b)+L(a)X(b)-X(ab),\\
d_RY(a,b)&=Y(b)R(a)+R(b)Y(a)-Y(ab),\\
\chi_{L,R}(X,Y)(a,b)&=[X(a),R(b)]+[L(a),Y(b)].
\end{align*}

\begin{definition}[Strict tangent observability operator]
Define
\[
\cD_p(X,Y)=
\begin{pmatrix}
\cO_B(X,Y)\\
X(1)\\
Y(1)\\
d_LX\\
d_RY\\
\chi_{L,R}(X,Y)
\end{pmatrix}.
\]
\end{definition}

\begin{theorem}[Strict tangent-space formula]\label{thm:tangent-formula}
The Zariski tangent space of the strict reconstruction scheme at $p$ is
\[
T_p\mathbf R_B=\ker\cD_p.
\]
\end{theorem}

\begin{proof}
Substitute $L_\varepsilon=L+\varepsilon X$ and $R_\varepsilon=R+\varepsilon Y$ into the weak Leibniz and bimodule equations over $\C[\varepsilon]/(\varepsilon^2)$.  The coefficient of $\varepsilon$ is exactly $\cD_p(X,Y)$.
\end{proof}

\begin{definition}[Three notions of rigidity]
A strict reconstruction $p$ is:
\begin{enumerate}[label=(\roman*)]
\item \emph{infinitesimally rigid} if $T_p\mathbf R_B=0$;
\item \emph{locally rigid} if $p$ is isolated in the reduced reconstruction variety;
\item \emph{globally unique} if $\Rfib(B)=\{p\}$.
\end{enumerate}
\end{definition}

\begin{proposition}[Infinitesimal rigidity implies local rigidity]\label{prop:inf-local}
In the finite-dimensional setting, if $T_p\mathbf R_B=0$, then $p$ is a reduced isolated point.
\end{proposition}

\begin{proof}
The local coordinate ring is Noetherian.  Vanishing tangent space gives $\mathfrak m_p/\mathfrak m_p^2=0$.  Nakayama's lemma yields $\mathfrak m_p=0$, so the local ring is $\C$.
\end{proof}

For the four points in Theorem~\ref{thm:four-points}, exact Jacobian computation gives full rank $32$ at every point.  Therefore each point is infinitesimally and locally rigid, while the reconstruction is not globally unique.  This separates local rigidity from global identifiability.

\section{Second-order obstruction theory}

Let $p=(L,R)\in\Rfib(B)$ and let
\[
\xi=(X,Y)\in T_p\mathbf R_B.
\]
We ask whether $\xi$ extends to a strict reconstruction over $\C[t]/(t^3)$:
\[
L_t=L+tX+t^2X_2,
\qquad
R_t=R+tY+t^2Y_2.
\]
The weak Leibniz equations are linear in $(L,R)$, so their second-order source term is zero.  The quadratic bimodule equations produce
\begin{align}
d_LX_2(a,b)&=-X(a)X(b), \label{eq:second-left}\\
d_RY_2(a,b)&=-Y(b)Y(a), \label{eq:second-right}\\
\chi_{L,R}(X_2,Y_2)(a,b)&=-[X(a),Y(b)]. \label{eq:second-comm}
\end{align}
Together with $\cO_B(X_2,Y_2)=0$ and $X_2(1)=Y_2(1)=0$, these equations can be written as
\[
\cD_p(X_2,Y_2)=-\cQ_p(X,Y),
\]
where
\[
\cQ_p(X,Y)=
\begin{pmatrix}
0\\0\\0\\
((a,b)\mapsto X(a)X(b))\\
((a,b)\mapsto Y(b)Y(a))\\
((a,b)\mapsto [X(a),Y(b)])
\end{pmatrix}.
\]

\begin{definition}[Second-order obstruction space and map]
Set
\[
\Obs_p(B)=\coker\cD_p
\]
and define
\[
\operatorname{ob}_{2,p}:T_p\mathbf R_B\longrightarrow\Obs_p(B),
\qquad
\operatorname{ob}_{2,p}(X,Y)=[\cQ_p(X,Y)].
\]
\end{definition}

\begin{theorem}[Primary obstruction criterion]\label{thm:obstruction}
A strict tangent vector $\xi\in T_p\mathbf R_B$ extends to order two if and only if
\[
\operatorname{ob}_{2,p}(\xi)=0.
\]
\end{theorem}

\begin{proof}
An order-two extension exists precisely when the affine equation
\[
\cD_p(X_2,Y_2)=-\cQ_p(\xi)
\]
has a solution.  This is equivalent to $\cQ_p(\xi)$ lying in $\operatorname{im}\cD_p$.
\end{proof}

The map $\operatorname{ob}_{2,p}$ is homogeneous quadratic rather than linear.  Its polarization is a symmetric bilinear expression whose three nonzero components are
\begin{align*}
X(a)X'(b)+X'(a)X(b),\\
Y(b)Y'(a)+Y'(b)Y(a),\\
[X(a),Y'(b)]+[X'(a),Y(b)].
\end{align*}
Section~\ref{sec:dgla} shows that this polarization is already the degree-one bracket of a two-step nilpotent differential graded Lie algebra; no higher $L_\infty$ operations are required for the present framed problem.

\section{A nonreduced false tangent}\label{sec:false-tangent}

Let
\[
A=\C[\epsilon]/(\epsilon^2),
\qquad
V=\C,
\]
and define $B:A\times A\to\C$ by
\[
B(\epsilon,\epsilon)=1
\]
and $B=0$ on the other basis pairs.  Then $V=\Span B(A,A)$.

A unital one-dimensional action pair is determined by
\[
L(\epsilon)=\ell,
\qquad
R(\epsilon)=r.
\]
The weak Leibniz equations give
\[
\ell+r=0,
\]
while the representation equations give
\[
\ell^2=0,
\qquad
r^2=0.
\]
Hence
\[
\mathbf R_B
=\Spec\frac{\C[\ell,r]}{(\ell+r,\ell^2,r^2)}
\cong
\Spec\frac{\C[\ell]}{(\ell^2)}.
\]
The scheme has one geometric point $p=(0,0)$ but a one-dimensional tangent space
\[
T_p\mathbf R_B=\C(1,-1).
\]

\begin{theorem}[A first-order direction obstructed at second order]\label{thm:false-tangent}
Every nonzero vector in $T_p\mathbf R_B$ is obstructed at second order.
\end{theorem}

\begin{proof}
For a tangent vector $(x,-x)$, attempt
\[
\ell_t=tx+t^2u,
\qquad
r_t=-tx+t^2v.
\]
The weak equation requires $u+v=0$.  However,
\[
\ell_t^2=t^2x^2\pmod{t^3},
\qquad
r_t^2=t^2x^2\pmod{t^3}.
\]
Over $\C$, these vanish only when $x=0$.
\end{proof}

Thus a nonzero Zariski tangent vector need not arise from any genuine local family.  It may record only the nilpotent thickness of a nonreduced reconstruction scheme.

\section{Higher-order recursion}

Suppose
\[
L_t=L+\sum_{n\ge1}t^nX_n,
\qquad
R_t=R+\sum_{n\ge1}t^nY_n.
\]
At order $n\ge2$, the extension equation has the form
\[
\cD_p(X_n,Y_n)=-\cQ_n,
\]
where the nonzero source terms are
\begin{align*}
\cQ_n^L(a,b)&=\sum_{i=1}^{n-1}X_i(a)X_{n-i}(b),\\
\cQ_n^R(a,b)&=\sum_{i=1}^{n-1}Y_i(b)Y_{n-i}(a),\\
\cQ_n^C(a,b)&=\sum_{i=1}^{n-1}[X_i(a),Y_{n-i}(b)].
\end{align*}
The class $[\cQ_n]\in\coker\cD_p$ is the order-$n$ obstruction relative to the previously chosen lifts.  Section~\ref{sec:dgla} packages the entire recursion into a Maurer--Cartan equation; the Artin-ring viewpoint remains the natural language for functorial lifting problems in the sense of \cite{Schlessinger1968}.


\section{Deformation DGLAs}\label{sec:dgla}

The preceding tangent and obstruction formulas are not merely analogous to deformation theory.  They arise from explicit differential graded Lie algebras.  We first record a finite presentation model that reproduces the complete polynomial system, and then a more intrinsic Hochschild model for framed deformations.

\subsection{A two-step presentation DGLA}

Fix $p=(L,R)\in\Rfib(B)$ and write $E=\End(V)$.  Let
\[
\gpres_{B,p}^{1}=\Hom(A,E)^2
\]
and let the relation space be
\[
\Rel_B=
\Hom(A^{\otimes3},V)^2
\oplus E^2
\oplus\Hom(A^{\otimes2},E)^3.
\]
Set
\[
\gpres_{B,p}^{2}=\Rel_B,
\qquad
\gpres_{B,p}^{n}=0\quad(n\neq1,2).
\]
The differential is
\[
d_p\big|_{\gpres^1}=\cD_p,
\qquad
d_p\big|_{\gpres^2}=0.
\]

For $\xi=(X,Y)$ and $\eta=(X',Y')$ in degree one, define $[\xi,\eta]_p\in\gpres^2$ to have zero Leibniz and unit components and the following three relation components:
\begin{align}
[\xi,\eta]^L_p(a,b)
&=X(a)X'(b)+X'(a)X(b),\\
[\xi,\eta]^R_p(a,b)
&=Y(b)Y'(a)+Y'(b)Y(a),\\
[\xi,\eta]^C_p(a,b)
&=[X(a),Y'(b)]+[X'(a),Y(b)].
\end{align}
Every bracket involving a degree-two element is defined to be zero.

\begin{proposition}[Presentation DGLA]\label{prop:presentation-dgla}
The graded vector space
\[
\gpres_{B,p}=\gpres_{B,p}^{1}\oplus\gpres_{B,p}^{2}
\]
with differential $d_p$ and bracket $[-,-]_p$ is a two-step nilpotent differential graded Lie algebra.
\end{proposition}

\begin{proof}
Since $d_p$ vanishes on degree two, $d_p^2=0$.  The bracket of two degree-one elements is symmetric, as required by graded skew-symmetry in odd degree.  Every nested bracket vanishes because degree two is central, so the graded Jacobi identity holds.  The derivation identity for $d_p$ is also immediate: both sides vanish after a bracket lands in degree two.
\end{proof}

\begin{theorem}[Exact Maurer--Cartan description]\label{thm:presentation-mc}
For $\xi=(X,Y)\in\gpres_{B,p}^{1}$,
\[
(L+X,R+Y)\in\Rfib(B)
\]
if and only if
\[
d_p\xi+\frac12[\xi,\xi]_p=0.
\]
Consequently,
\[
H^1(\gpres_{B,p})=T_p\mathbf R_B,
\qquad
H^2(\gpres_{B,p})=\coker\cD_p=\Obs_p(B),
\]
and
\[
\operatorname{ob}_{2,p}(\xi)
=
\left[\frac12[\xi,\xi]_p\right].
\]
\end{theorem}

\begin{proof}
The weak Leibniz and unit equations are affine-linear and contribute only $d_p\xi$.  Expanding the left representation, right antirepresentation, and commuting-action equations at $(L,R)$ gives respectively
\[
(a,b)\longmapsto d_LX(a,b)+X(a)X(b),
\]
\[
(a,b)\longmapsto d_RY(a,b)+Y(b)Y(a),
\]
and
\[
(a,b)\longmapsto\chi_{L,R}(X,Y)(a,b)+[X(a),Y(b)],
\]
which are exactly the corresponding components of $d_p\xi+\tfrac12[\xi,\xi]_p$.  The cohomology identities follow because $\gpres^0=\gpres^3=0$.
\end{proof}

\begin{remark}[Presentation dependence]
The space $H^2(\gpres_{B,p})$ is a correct recipient for obstruction classes, but it need not be minimal.  The chosen list of defining equations can contain linear dependencies and higher syzygies, which enlarge the cokernel without producing additional geometric obstructions.  This motivates the intrinsic model below.
\end{remark}

\subsection{A constrained Hochschild DGLA}

Let
\[
S=A^e=A\otimes A^{\mathrm{op}}.
\]
The bimodule structure $(L,R)$ is equivalent to the algebra homomorphism
\[
\mu:S\longrightarrow E,
\qquad
\mu(a\otimes b^{\mathrm{op}})=L(a)R(b).
\]
Regard $E$ as an $S$-bimodule through $\mu$, and let
\[
C^n(S,E_\mu)=\Hom(\overline{S}^{\otimes n},E),
\qquad
\overline S=S/\C1,
\]
be the normalized Hochschild cochains.  Their differential is
\begin{align*}
(d_\mu f)(s_1,\ldots,s_{n+1})
={}&\mu(s_1)f(s_2,\ldots,s_{n+1})\\
&+\sum_{i=1}^{n}(-1)^if(s_1,\ldots,s_is_{i+1},\ldots,s_{n+1})\\
&+(-1)^{n+1}f(s_1,\ldots,s_n)\mu(s_{n+1}).
\end{align*}
The cup product is
\[
(f\smile g)(s_1,\ldots,s_{m+n})
=f(s_1,\ldots,s_m)g(s_{m+1},\ldots,s_{m+n}),
\]
and its graded commutator makes the Hochschild differential graded algebra into a DGLA.

For $f\in C^1(S,E_\mu)$, put
\[
X_f(a)=f(a\otimes1),
\qquad
Y_f(b)=f(1\otimes b^{\mathrm{op}}),
\]
and define the linear $B$-compatibility operator
\[
\lambda_B:C^1(S,E_\mu)\longrightarrow\Hom(A^{\otimes3},V)^2
\]
by
\begin{align*}
\lambda_B(f)_1(a,b,c)
&=X_f(a)B(b,c)+Y_f(b)B(a,c),\\
\lambda_B(f)_2(a,b,c)
&=Y_f(c)B(a,b)+X_f(b)B(a,c).
\end{align*}
Define
\[
\ghoch_{B,p}^{1}=\ker\lambda_B,
\qquad
\ghoch_{B,p}^{n}=C^n(S,E_\mu)\quad(n\ge2),
\qquad
\ghoch_{B,p}^{n}=0\quad(n\le0).
\]

\begin{proposition}[Constrained Hochschild DGLA]\label{prop:hoch-dgla}
The graded subspace $\ghoch_{B,p}$ is a DGLA.  For every local Artin $\C$-algebra $(R_{\mathrm{Art}},\mathfrak m)$, its Maurer--Cartan elements in $\ghoch_{B,p}\otimes\mathfrak m$ are in bijection with framed strict reconstruction deformations of $(L,R)$ over $R_{\mathrm{Art}}$ that keep $A$, $V$, and $B$ fixed.
\end{proposition}

\begin{proof}
The only additional condition occurs in degree one.  Its differential and the bracket of two degree-one elements land in the unrestricted degree-two space, so $\ghoch_{B,p}$ is closed under both operations.  For $\alpha\in C^1(S,E_\mu)\otimes\mathfrak m$, the map $\mu+\alpha$ is an algebra homomorphism precisely when
\[
d_\mu\alpha+\alpha\smile\alpha
=d_\mu\alpha+\frac12[\alpha,\alpha]=0.
\]
Because the Leibniz equations are linear in the actions, preservation of $B$ is exactly the condition $\lambda_B(\alpha)=0$.
\end{proof}

\begin{theorem}[Intrinsic tangent and primary obstruction]\label{thm:intrinsic-obstruction}
There is a natural isomorphism
\[
H^1(\ghoch_{B,p})\cong T_p\mathbf R_B.
\]
If $\xi=(X,Y)$ corresponds to the cocycle $f_\xi$, then its intrinsic primary obstruction is
\[
\operatorname{Ob}^{\mathrm{Hoch}}_{2,p}(\xi)
=[f_\xi\smile f_\xi]
\in H^2(\ghoch_{B,p}),
\]
and $\xi$ extends to second order if and only if this class vanishes.
\end{theorem}

\begin{proof}
If $d_\mu f=0$, then
\[
f(a\otimes b^{\mathrm{op}})
=X_f(a)R(b)+L(a)Y_f(b),
\]
so a normalized Hochschild cocycle is determined by its two restrictions.  The cocycle equation on the two factor subalgebras and on mixed products yields precisely the linearized left representation, right antirepresentation, and commuting-action equations; $\lambda_B(f)=0$ gives the two linearized Leibniz equations.  Conversely, a strict tangent pair defines a cocycle by the displayed formula.  Since degree zero has been removed in the framed problem, cocycles are not quotiented by infinitesimal conjugations.

For a first-order cocycle $f_\xi$, the coefficient of $t^2$ in the Maurer--Cartan equation is
\[
d_\mu f_2=-f_\xi\smile f_\xi.
\]
The required correction must lie in $\ghoch^1_{B,p}$, so solvability is equivalent to vanishing of the stated class in $H^2(\ghoch_{B,p})$.
\end{proof}

\subsection{A relation-level DGLA morphism}

For $f\in\ghoch_{B,p}^{1}$, define
\[
\Theta(f)=(X_f,Y_f)\in\gpres_{B,p}^{1}.
\]
For a Hochschild two-cochain $\phi$, define
\[
\Pi(\phi)=
\begin{pmatrix}
(a,b)\mapsto\phi(a\otimes1,b\otimes1)\\
(a,b)\mapsto\phi(1\otimes b^{\mathrm{op}},1\otimes a^{\mathrm{op}})\\
(a,b)\mapsto\phi(a\otimes1,1\otimes b^{\mathrm{op}})-\phi(1\otimes b^{\mathrm{op}},a\otimes1)
\end{pmatrix}.
\]
Let
\[
\iota_{\mathrm{rel}}:
\Hom(A^{\otimes2},E)^3\longrightarrow\Rel_B
\]
insert zero weak-Leibniz and unit components.  Define graded maps
\[
\Psi^1=\Theta,
\qquad
\Psi^2=\iota_{\mathrm{rel}}\circ\Pi,
\qquad
\Psi^n=0\quad(n\ge3).
\]

\begin{proposition}[Relation DGLA morphism]\label{prop:relation-dgla-morphism}
The maps $\Psi^n$ define a morphism of DGLAs
\[
\Psi:\ghoch_{B,p}\longrightarrow\gpres_{B,p}.
\]
\end{proposition}

\begin{proof}
For $f\in\ghoch^1_{B,p}$, normalizedness gives $X_f(1)=Y_f(1)=0$, while $\lambda_B(f)=0$ gives the two linearized weak-Leibniz components.  On the two factor subalgebras and on mixed products, direct expansion gives
\begin{align*}
\Pi(d_\mu f)_L(a,b)&=d_LX_f(a,b),\\
\Pi(d_\mu f)_R(a,b)&=d_RY_f(a,b),\\
\Pi(d_\mu f)_C(a,b)&=\chi_{L,R}(X_f,Y_f)(a,b).
\end{align*}
Hence
\[
\Psi^2(d_\mu f)=d_p\Psi^1(f).
\]
The target differential vanishes on degree two and $\Psi$ vanishes in degrees at least three, so this proves compatibility with the differential in all degrees.

For $f,g\in\ghoch^1_{B,p}$, the degree-one Hochschild bracket is
\[
[f,g]=f\smile g+g\smile f.
\]
Applying $\Pi$ gives
\begin{align*}
\Pi([f,g])_L(a,b)
&=X_f(a)X_g(b)+X_g(a)X_f(b),\\
\Pi([f,g])_R(a,b)
&=Y_f(b)Y_g(a)+Y_g(b)Y_f(a),\\
\Pi([f,g])_C(a,b)
&=[X_f(a),Y_g(b)]+[X_g(a),Y_f(b)],
\end{align*}
which is precisely $[\Theta(f),\Theta(g)]_p$.  Every bracket involving degree two maps to zero on both sides.  Thus $\Psi$ is a DGLA morphism.
\end{proof}

\begin{corollary}[Cohomological obstruction projection]\label{cor:cohomological-projection}
The DGLA morphism $\Psi$ induces natural maps
\[
H^1(\Psi):H^1(\ghoch_{B,p})\longrightarrow H^1(\gpres_{B,p}),
\qquad
\overline\Pi:=H^2(\Psi):H^2(\ghoch_{B,p})\longrightarrow\coker\cD_p.
\]
Under the tangent-space identifications, $H^1(\Psi)$ is the identity on $T_p\mathbf R_B$.  Moreover, for every $\xi\in T_p\mathbf R_B$,
\[
\boxed{
\overline\Pi\bigl(\operatorname{Ob}^{\mathrm{Hoch}}_{2,p}(\xi)\bigr)
=
\operatorname{ob}_{2,p}(\xi).
}
\]
Equivalently, the diagram
\[
\begin{CD}
T_p\mathbf R_B @>{\operatorname{Ob}^{\mathrm{Hoch}}_{2,p}}>> H^2(\ghoch_{B,p})\\
@V{\mathrm{id}}VV @VV{\overline\Pi}V\\
T_p\mathbf R_B @>{\operatorname{ob}_{2,p}}>> \coker\cD_p
\end{CD}
\]
commutes.
\end{corollary}

\begin{proof}
The first assertion follows from Theorems~\ref{thm:presentation-mc} and \ref{thm:intrinsic-obstruction} together with the degree-one definition of $\Psi$.  If $f_\xi$ represents $\xi=(X,Y)$, then
\[
\Pi(f_\xi\smile f_\xi)
=
\begin{pmatrix}
(a,b)\mapsto X(a)X(b)\\
(a,b)\mapsto Y(b)Y(a)\\
(a,b)\mapsto [X(a),Y(b)]
\end{pmatrix},
\]
which is exactly the nonzero relation part of $\cQ_p(\xi)$.  Passing to cohomology gives the claimed identity.
\end{proof}

\subsection{Exact cohomology and the minimal Kuranishi sector for the dual-number example}

For the example of Section~\ref{sec:false-tangent}, write
\[
x=\epsilon\otimes1,
\qquad
y=1\otimes\epsilon^{\mathrm{op}},
\qquad
w=xy.
\]
Then $\overline S$ has basis $x,y,w$, with the only nonzero products in $\overline S$ given by
\[
xy=yx=w.
\]
Let $\delta_{uv}$ denote the normalized two-cochain that takes value $1$ on $(u,v)$ and vanishes on the other basis pairs.

The $B$-compatibility condition on a normalized one-cochain is
\[
f(x)+f(y)=0.
\]
Thus $\ghoch^1_{B,p}$ has basis
\[
\tau=(-1,1,0),
\qquad
u=(0,0,1),
\]
where the triples record the values on $(x,y,w)$.  The Hochschild differential satisfies
\[
d_\mu\tau=0,
\qquad
d_\mu\nu=\beta:=-(\delta_{xy}+\delta_{yx}).
\]
Consequently,
\[
H^1(\ghoch_{B,p})=\C[\tau].
\]

Exact calculation gives
\[
Z^2(\ghoch_{B,p})
=
\Span\{\eta,\eta_1,\eta_2,\beta\},
\]
where
\begin{align*}
\eta&=\tau\smile\tau
=\delta_{xx}-\delta_{xy}-\delta_{yx}+\delta_{yy},\\
\eta_1&=\delta_{xx}-\delta_{yy},\\
\eta_2&=\delta_{xy}-\delta_{yx}.
\end{align*}
Since $B^2=\C\beta$, one obtains
\[
H^2(\ghoch_{B,p})
=
\C[\eta]\oplus\C[\eta_1]\oplus\C[\eta_2].
\]
In particular,
\[
\dim\ghoch^1_{B,p}=2,
\qquad
\dim Z^1=1,
\qquad
\dim Z^2=4,
\qquad
\dim H^2=3.
\]

\begin{proposition}[The intrinsic relation map is injective in degree two]\label{prop:dual-h2-injective}
For the dual-number example, the map
\[
\overline\Pi:H^2(\ghoch_{B,p})\longrightarrow\coker\cD_p
\]
is injective.
\end{proposition}

\begin{proof}
Restrict to the three left, right, and mixed relation coordinates evaluated at $(\epsilon,\epsilon)$.  The three cohomology generators have images
\[
\overline\Pi[\eta]=(1,1,0),
\qquad
\overline\Pi[\eta_1]=(1,-1,0),
\qquad
\overline\Pi[\eta_2]=(0,0,2).
\]
At the origin $L(\epsilon)=R(\epsilon)=0$, and $\epsilon^2=0$, so the linearized left and right representation equations vanish identically in these two $(\epsilon,\epsilon)$ coordinates.  The mixed coordinate also vanishes on $\operatorname{im}\cD_p$ because $E=\End(\C)=\C$ is commutative.  The three displayed vectors therefore remain linearly independent in the cokernel.
\end{proof}

\begin{theorem}[Minimal Kuranishi equation for the false tangent]\label{thm:dual-minimal-kuranishi}
There is a homotopy transfer of $\ghoch_{B,p}$ to its cohomology for which the operations on the $H^1$-generated Maurer--Cartan sector satisfy
\[
\ell_2([\tau],[\tau])=2[\eta],
\qquad
\ell_n([\tau],\ldots,[\tau])=0
\quad(n\ge3).
\]
Hence the minimal Kuranishi map is
\[
\kappa(t[\tau])=t^2[\eta].
\]
For every local Artin algebra $(R_{\mathrm{Art}},\mathfrak m)$, the framed deformation functor of the dual-number reconstruction point is therefore
\[
\MC(R_{\mathrm{Art}})
\cong
\{t\in\mathfrak m\mid t^2=0\},
\]
and its prorepresenting hull is
\[
\boxed{
\C[[t]]/(t^2).
}
\]
This agrees with the completed local ring of
\[
\mathbf R_B\cong\Spec\C[\ell]/(\ell^2)
\]
at its unique geometric point.
\end{theorem}

\begin{proof}
Choose representatives of $H^2$ to be $\eta,\eta_1,\eta_2$, and choose a degree $-1$ homotopy with
\[
h(\beta)=\nu,
\qquad
h(\eta)=h(\eta_1)=h(\eta_2)=0.
\]
Since all spaces are vector spaces over $\C$, this low-degree splitting extends degreewise to a contraction of the full cochain complex onto its cohomology.  The transferred binary bracket is the projection of the original DGLA bracket, so
\[
\ell_2([\tau],[\tau])
=\pi[\tau,\tau]
=2[\tau\smile\tau]
=2[\eta].
\]
Moreover,
\[
h[\tau,\tau]=2h(\eta)=0.
\]
Every rooted binary tree contributing to $\ell_n([\tau],\ldots,[\tau])$ for $n\ge3$ contains a lowest internal edge carrying the factor $h[\tau,\tau]$; hence every such tree vanishes.  The minimal Maurer--Cartan equation is therefore
\[
\frac12\ell_2(t[\tau],t[\tau])=t^2[\eta]=0.
\]
Since $[\eta]\neq0$, this is equivalent to $t^2=0$.  The final identification follows from the explicit strict reconstruction equation of Section~\ref{sec:false-tangent}.
\end{proof}

\begin{remark}[Cohomological receptacle versus effective obstruction image]
The equality $\dim H^2=3$ does not mean that the one-dimensional local germ requires three equations.  Proposition~\ref{prop:dual-h2-injective} shows that all three classes are genuine independent relation-level classes, not projection artifacts.  Nevertheless, the Kuranishi map from the one-dimensional tangent space has image
\[
\operatorname{im}\kappa=\C[\eta].
\]
Thus one must distinguish the full obstruction receptacle $H^2$ from the effective obstruction directions actually reached by the Maurer--Cartan equation near a specified point.
\end{remark}

\begin{remark}[When higher $L_\infty$ operations become necessary]
For the dual-number point, Theorem~\ref{thm:dual-minimal-kuranishi} shows that all higher operations vanish on the $H^1$-generated Maurer--Cartan sector.  In larger reconstruction problems, nonzero higher brackets may appear after homotopy transfer, taking a homotopy fiber or gauge quotient, or allowing the multiplication of $A$ and the map $B$ to deform simultaneously.
\end{remark}


\section{Effective obstruction spaces and minimal equations}\label{sec:effective-obstruction}

The preceding examples show that $H^2$ is generally too large to count the equations of the local reconstruction germ.  This section isolates the intrinsic finite-dimensional quotient that does count them.

Assume that $T=H^1(\ghoch_{B,p})$ and $O=H^2(\ghoch_{B,p})$ are finite-dimensional, and choose a transferred minimal $L_\infty$ structure on cohomology.  Its Maurer--Cartan equation on degree one is the formal Kuranishi map
\[
\kappa(\xi)
=
\sum_{n\ge2}\frac1{n!}\ell_n(\xi,\ldots,\xi)
\in O\widehat\otimes \widehat{\operatorname{Sym}}(T^\vee).
\]
Put
\[
P_T=\widehat{\operatorname{Sym}}(T^\vee),
\qquad
\mathfrak n=P_T^{\ge1},
\]
and let
\[
\kappa^\sharp:O^\vee\longrightarrow\mathfrak n^2,
\qquad
\lambda\longmapsto \lambda\circ\kappa.
\]
The \emph{Kuranishi ideal} is
\[
I_\kappa=P_T\,\kappa^\sharp(O^\vee).
\]
The framed formal reconstruction germ is represented by
\[
R_p=P_T/I_\kappa.
\]
Because the transferred structure is minimal, $I_\kappa\subseteq\mathfrak n^2$, so this is a minimal formal embedding: the cotangent space of $R_p$ is $T^\vee$.

\begin{definition}[Effective relation and obstruction spaces]\label{def:effective-obstruction}
The effective relation space and effective obstruction space at $p$ are
\[
\mathscr R^{\mathrm{eff}}_p
:=
I_\kappa/\mathfrak n I_\kappa,
\qquad
\Obs^{\mathrm{eff}}_p
:=
\bigl(\mathscr R^{\mathrm{eff}}_p\bigr)^\vee.
\]
\end{definition}

\begin{proposition}[Minimal-equation theorem]\label{prop:minimal-equations}
The number
\[
\dim_\C\Obs^{\mathrm{eff}}_p
=
\dim_\C I_\kappa/\mathfrak n I_\kappa
\]
is the minimum number of formal equations required to generate the defining ideal of the framed reconstruction germ inside a smooth formal space of minimal embedding dimension $\dim T$.
\end{proposition}

\begin{proof}
A family $f_1,\ldots,f_r\in I_\kappa$ generates $I_\kappa$ if and only if its residue classes span $I_\kappa/\mathfrak n I_\kappa$.  This is the complete-local form of Nakayama's lemma.  Thus every basis of $I_\kappa/\mathfrak n I_\kappa$ lifts to a minimal generating family, and no smaller family can generate the ideal.
\end{proof}

\begin{theorem}[Independence of the minimal Kuranishi presentation]\label{thm:effective-intrinsic}
Let $R$ be a complete Noetherian local $\C$-algebra with residue field $\C$.  Suppose
\[
P=\C[[x_1,\ldots,x_e]]\twoheadrightarrow R,
\qquad
P'=\C[[y_1,\ldots,y_e]]\twoheadrightarrow R
\]
are two minimal presentations, with kernels $I\subseteq\mathfrak n^2$ and $I'\subseteq(\mathfrak n')^2$.  Then a formal change of coordinates identifies the presentations and induces an isomorphism
\[
I/\mathfrak n I
\cong
I'/\mathfrak n'I'.
\]
Consequently, the isomorphism class of $\mathscr R^{\mathrm{eff}}_p$, the isomorphism class of $\Obs^{\mathrm{eff}}_p$, and their common dimension are independent of the chosen contraction and minimal Kuranishi presentation.
\end{theorem}

\begin{proof}
Both sets of variables map to minimal generating systems of the maximal ideal of $R$, hence their images give bases of $\mathfrak m_R/\mathfrak m_R^2$.  Lift the images of the $x_i$ to formal series in the $y_j$.  This defines a continuous homomorphism
\[
\varphi:P\longrightarrow P'
\]
commuting with the maps to $R$.  Its linear part is the change-of-basis matrix between two bases of $\mathfrak m_R/\mathfrak m_R^2$, hence is invertible.  The formal inverse function theorem makes $\varphi$ an isomorphism.  Since the two maps to $R$ commute with $\varphi$, one has $\varphi(I)=I'$, and therefore
\[
I/\mathfrak n I\xrightarrow{\sim}I'/\mathfrak n'I'.
\]
Applying this to two minimal Kuranishi presentations proves the final assertion.
\end{proof}

\begin{remark}[Cotangent-complex interpretation]
The preceding relation space is the minimal-presentation realization of the first relation homology of the local germ.  In cotangent-complex language, it is the completed local analogue of
\[
I/\mathfrak n I
\cong
H_1\bigl(L_{R/\C}\otimes_R^{\mathbf L}\C\bigr)
\]
for a minimal regular presentation.  The vanishing of the linear part of the relations makes the differential to the cotangent space zero; compare the standard presentation description of the cotangent complex in \cite{StacksCotangent}.
\end{remark}

\begin{proposition}[Embedding into the cohomological obstruction receptacle]\label{prop:effective-into-h2}
The component map of the Kuranishi equation induces a natural surjection
\[
q_\kappa:O^\vee\twoheadrightarrow\mathscr R^{\mathrm{eff}}_p,
\qquad
\lambda\longmapsto[\lambda\circ\kappa].
\]
Hence, when $O$ is finite-dimensional, duality gives an injection
\[
j_\kappa:\Obs^{\mathrm{eff}}_p\hookrightarrow O=H^2(\ghoch_{B,p}).
\]
The abstract effective obstruction space is intrinsic by Theorem~\ref{thm:effective-intrinsic}; its displayed embedding into a fixed coordinate realization of $H^2$ is determined only up to the linear isomorphism induced by a change of minimal model.
\end{proposition}

\begin{proof}
The ideal $I_\kappa$ is generated by the image of $\kappa^\sharp$, so the residue classes of those components span $I_\kappa/\mathfrak nI_\kappa$.  This proves surjectivity.  Dualizing a surjection of finite-dimensional vector spaces gives the injection $j_\kappa$.
\end{proof}

Combining Proposition~\ref{prop:effective-into-h2} with Corollary~\ref{cor:cohomological-projection} gives the three-stage comparison
\[
\boxed{
\Obs^{\mathrm{eff}}_p
\xhookrightarrow{\ j_\kappa\ }
H^2(\ghoch_{B,p})
\xrightarrow{\ \overline\Pi\ }
\coker\cD_p.
}
\]
We call
\[
\Obs^{\mathrm{rel,eff}}_p
:=
\operatorname{im}(\overline\Pi\circ j_\kappa)
\]
the effective relation-level obstruction space.  The map
\[
\overline\Pi_{\mathrm{eff}}
:=
\overline\Pi\circ j_\kappa
\]
need not be injective in general; its kernel and the complement of its image in $\coker\cD_p$ measure different failures and should not both be called presentation syzygies.

\begin{corollary}[Effective obstruction for the dual-number germ]\label{cor:dual-effective}
For the dual-number example,
\[
I_\kappa=(t^2),
\qquad
\mathscr R^{\mathrm{eff}}_p
=(t^2)/(t^3),
\qquad
\dim\Obs^{\mathrm{eff}}_p=1.
\]
Under the basis of Section~\ref{sec:dgla}, the injection $j_\kappa$ has image
\[
\operatorname{im}j_\kappa=\C[\eta]\subset H^2(\ghoch_{B,p}).
\]
Moreover, $\overline\Pi_{\mathrm{eff}}$ is injective.
\end{corollary}

\begin{proof}
The first three assertions follow immediately from $I_\kappa=(t^2)$.  Since
\[
\kappa(t[\tau])=t^2[\eta],
\]
the dual component map kills $\eta_1^\vee$ and $\eta_2^\vee$ and maps $\eta^\vee$ to the generator of $(t^2)/(t^3)$.  Thus the dual injection lands on $\C[\eta]$.  Proposition~\ref{prop:dual-h2-injective} then implies injectivity after applying $\overline\Pi$.
\end{proof}

\subsection{A two-parameter mixed-obstruction example}\label{sec:two-parameter-example}

Let
\[
A=\C1\oplus\C a\oplus\C b,
\qquad
(a,b)^2=0,
\]
and let $V=\C^2$ with basis $e_1,e_2$.  Define
\[
B(b,a)=e_1,
\qquad
B(b,b)=e_2,
\]
and let all other values on basis pairs vanish.  Then $V=\Span B(A,A)$.  Let $p$ be the augmentation bimodule structure
\[
L(x)=R(x)=\varepsilon(x)I_V.
\]
Because $B$ vanishes whenever one input is scalar and the radical of $A$ has square zero, $B$ is a biderivation relative to $p$.

\begin{proposition}[Exact strict reconstruction scheme]\label{prop:two-parameter-scheme}
Every unit-preserving weak reconstruction has
\[
L(a)=R(a)=0,
\qquad
L(b)=N(u,v),
\qquad
R(b)=-N(u,v),
\]
where
\[
N(u,v)=
\begin{pmatrix}
0&-u\\
0&-v
\end{pmatrix}.
\]
It is strict precisely when
\[
uv=0,
\qquad
v^2=0.
\]
Consequently,
\[
\boxed{
\mathbf R_B
\cong
\Spec\frac{\C[u,v]}{(uv,v^2)}.
}
\]
\end{proposition}

\begin{proof}
Evaluating the two weak Leibniz identities on the eight triples from $\{a,b\}^3$ gives a linear system in the sixteen matrix entries of
\[
L(a),L(b),R(a),R(b).
\]
Exact row reduction has rank fourteen and yields precisely the displayed two-parameter solution.  Since the radical of $A$ has square zero, all left and right representation equations reduce to $N(u,v)^2=0$; the mixed equations are automatic because $R(b)=-L(b)$.  Finally,
\[
N(u,v)^2
=
\begin{pmatrix}
0&uv\\
0&v^2
\end{pmatrix},
\]
which gives the stated ideal.
\end{proof}

\begin{remark}[A smooth branch with an embedded false direction]
The reduced reconstruction scheme is the affine line $v=0$.  The $u$-direction integrates to an actual family of strict reconstructions, while the $v$-direction exists in the Zariski tangent space only at the origin and is nilpotently embedded there.  The mixed equation $uv=0$ records the failure to combine the genuine $u$-deformation with the false $v$-direction.
\end{remark}

For the intrinsic complex, put
\[
x=b\otimes1,
\qquad
y=1\otimes b^{\mathrm{op}},
\]
and define matrices
\[
P=
\begin{pmatrix}0&-1\\0&0\end{pmatrix},
\qquad
Q=
\begin{pmatrix}0&0\\0&-1\end{pmatrix}.
\]
Define normalized one-cochains $\alpha,\beta$ by
\[
\alpha(x)=P,
\quad
\alpha(y)=-P,
\qquad
\beta(x)=Q,
\quad
\beta(y)=-Q,
\]
and let them vanish on the other normalized basis elements of $A^e$.

\begin{theorem}[Two-parameter minimal Kuranishi equation]\label{thm:two-parameter-kuranishi}
For the square-zero example,
\[
\dim\ghoch^1_{B,p}=18,
\qquad
\dim H^1(\ghoch_{B,p})=2,
\]
\[
\dim Z^2(\ghoch_{B,p})=64,
\qquad
\dim H^2(\ghoch_{B,p})=48.
\]
The classes $[\alpha],[\beta]$ form a basis of $H^1$.  Put
\[
\eta_{uv}=\alpha\smile\beta+\beta\smile\alpha,
\qquad
\eta_{vv}=\beta\smile\beta.
\]
Then $[\eta_{uv}]$ and $[\eta_{vv}]$ are nonzero and linearly independent in $H^2$.  There is a homotopy transfer for which
\begin{align*}
\ell_2([\alpha],[\alpha])&=0,\\
\ell_2([\alpha],[\beta])&=[\eta_{uv}],\\
\ell_2([\beta],[\beta])&=2[\eta_{vv}],
\end{align*}
and every $\ell_n$ with $n\ge3$ vanishes on the $H^1$-generated Maurer--Cartan sector.  Hence
\[
\boxed{
\kappa(u[\alpha]+v[\beta])
=
uv[\eta_{uv}]+v^2[\eta_{vv}].
}
\]
The Kuranishi ideal is $(uv,v^2)$, in agreement with Proposition~\ref{prop:two-parameter-scheme}.
\end{theorem}

\begin{proof}
The normalized augmentation ideal of $S=A^e$ has dimension eight.  Exact rational linear algebra gives
\[
\rank\lambda_B=14,
\qquad
\rank(d_1|_{\ghoch^1_{B,p}})=16,
\qquad
\rank d_2=192,
\]
which yields the four displayed dimensions.  Direct substitution shows that $\alpha$ and $\beta$ are $B$-compatible cocycles and span $Z^1=H^1$.

One has
\[
\alpha\smile\alpha=0.
\]
The two cochains $\eta_{uv}$ and $\eta_{vv}$ are closed, each increases the boundary rank by one, and together increase it by two; hence their classes are nonzero and independent.  Choose them as cohomology representatives and choose the contracting homotopy to vanish on both.  Every higher transfer tree on $[\alpha],[\beta]$ then has a lowest internal edge carrying
\[
h[\alpha,\alpha],
\quad
h[\alpha,\beta],
\quad\text{or}\quad
h[\beta,\beta],
\]
all of which vanish.  Thus only the binary bracket contributes, and
\[
\frac12\ell_2(u[\alpha]+v[\beta],u[\alpha]+v[\beta])
=
uv[\eta_{uv}]+v^2[\eta_{vv}].
\]
The agreement with the exact reconstruction scheme completes the proof.
\end{proof}

\begin{corollary}[Forty-eight-dimensional receptacle, two effective equations]\label{cor:two-parameter-effective}
For the square-zero example,
\[
\dim H^2(\ghoch_{B,p})=48,
\qquad
\dim\Obs^{\mathrm{eff}}_p=2,
\]
and
\[
\operatorname{im}j_\kappa
=
\Span\{[\eta_{uv}],[\eta_{vv}]\}.
\]
The effective relation map $\overline\Pi_{\mathrm{eff}}$ is injective.
\end{corollary}

\begin{proof}
Since $I_\kappa=(uv,v^2)$ and $\mathfrak n=(u,v)$, the classes of $uv$ and $v^2$ form a basis of
\[
I_\kappa/\mathfrak nI_\kappa.
\]
The Kuranishi component map identifies the dual effective space with the span of the two stated cohomology classes.  Under the relation projection, their left representation components at $(b,b)$ are respectively
\[
PQ+QP=
\begin{pmatrix}0&1\\0&0\end{pmatrix},
\qquad
Q^2=
\begin{pmatrix}0&0\\0&1\end{pmatrix}.
\]
The linearized left relation at $(b,b)$ vanishes identically because $b^2=0$ and $L(b)=0$ at the base point.  These two matrix coordinates therefore remain independent in $\coker\cD_p$.
\end{proof}


\section{A first genuinely cubic Kuranishi obstruction}\label{sec:cubic-example}

The preceding nonreduced examples are quadratically controlled.  We now give an explicit inverse-Leibniz reconstruction in which the complete quadratic Kuranishi map on the strict tangent space vanishes, but the cubic term does not.

Let
\[
A=\C[\epsilon]/(\epsilon^4),
\qquad
V=\C^3,
\]
and let $v_1,v_2,v_3$ be the standard basis of $V$.  Put
\[
N=
\begin{pmatrix}
0&1&0\\
0&0&1\\
0&0&0
\end{pmatrix}.
\]
Define a strict base reconstruction $p=(L,R)$ by
\[
L(\epsilon)=N,
\qquad
R(\epsilon)=N+N^2,
\]
with the values on higher powers determined multiplicatively.  Thus
\[
L(\epsilon^2)=R(\epsilon^2)=N^2,
\qquad
L(\epsilon^3)=R(\epsilon^3)=0.
\]
The two actions commute because they are polynomials in $N$.

Define $B:A\times A\to V$ by declaring all values with a unit input to vanish and taking the following as its only nonzero values on radical basis pairs:
\begin{align*}
B(\epsilon,\epsilon)&=3v_1+v_2+v_3,\\
B(\epsilon,\epsilon^2)=B(\epsilon^2,\epsilon)&=3v_1+2v_2,\\
B(\epsilon,\epsilon^3)=B(\epsilon^3,\epsilon)&=3v_1,\\
B(\epsilon^2,\epsilon^2)&=4v_1.
\end{align*}
Direct substitution verifies both Leibniz identities relative to $p$.  Moreover,
\[
V=\Span B(A,A).
\]

Let $\mathcal K_B$ denote the vector space of unit-preserving weak perturbations at $p$, and let
\[
\cD_p\big|_{\mathcal K_B}:\mathcal K_B\longrightarrow\Rel_B
\]
be the strict tangent operator restricted to that weak space.

\begin{proposition}[Exact weak and strict tangent dimensions]\label{prop:cubic-tangent}
For the preceding reconstruction,
\[
\dim\mathcal K_B=9,
\qquad
\rank\bigl(\cD_p|_{\mathcal K_B}\bigr)=7,
\qquad
\dim T_p\mathbf R_B=2.
\]
A basis of the strict tangent space is given by $\alpha=(X_\alpha,Y_\alpha)$ and $\xi=(X_\xi,Y_\xi)$, where
\begin{align*}
X_\alpha(\epsilon)&=-N^2,
&Y_\alpha(\epsilon)&=N^2,
\end{align*}
and all their values on $\epsilon^2,\epsilon^3$ vanish, while
\begin{align*}
X_\xi(\epsilon)&=-N,
&X_\xi(\epsilon^2)&=-2N^2,
&X_\xi(\epsilon^3)&=0,\\
Y_\xi(\epsilon)&=N,
&Y_\xi(\epsilon^2)&=2N^2,
&Y_\xi(\epsilon^3)&=0.
\end{align*}
\end{proposition}

\begin{proof}
The two linearized Leibniz blocks form an exact rational matrix with fifty-four columns.  Its rank is forty-five, giving the nine-dimensional weak perturbation space.  Restricting the strict relation differential to this kernel gives a matrix of rank seven.  The two displayed perturbations lie in both kernels and are linearly independent, hence form a tangent basis.
\end{proof}

Write
\[
\cP_p(\rho,\sigma)
:=
\cQ_p(\rho+\sigma)-\cQ_p(\rho)-\cQ_p(\sigma)
\]
for the symmetric polarization of the quadratic relation map.  Define a weak second-order correction $\eta=(U,W)$ by
\begin{align*}
U(\epsilon)&=I,
&U(\epsilon^2)&=2N+N^2,
&U(\epsilon^3)&=3N^2,\\
W(\epsilon)&=-I,
&W(\epsilon^2)&=-2N-N^2,
&W(\epsilon^3)&=-3N^2.
\end{align*}
Then exact substitution gives
\begin{equation}\label{eq:cubic-second-lift}
\cQ_p(\alpha)=0,
\qquad
\cP_p(\alpha,\xi)=0,
\qquad
\cD_p\eta=-\cQ_p(\xi).
\end{equation}
Thus every quadratic Kuranishi product on $T_p\mathbf R_B$ vanishes in the obstruction cokernel.

At third order define
\[
\zeta:=\cP_p(\xi,\eta)\in\Rel_B.
\]
For a relation vector $\Theta$, let $\Theta^L(a,b)$ and $\Theta^R(a,b)$ denote its left and right representation components.  Consider the linear functional
\begin{equation}\label{eq:cubic-certificate}
\Lambda(\Theta)
=
4\,\Theta^L(\epsilon,\epsilon)_{12}
-2\,\Theta^L(\epsilon,\epsilon)_{23}
-2\,\Theta^L(\epsilon,\epsilon^2)_{13}
-2\,\Theta^R(\epsilon,\epsilon)_{23}.
\end{equation}

\begin{theorem}[Nonzero cubic Kuranishi term]\label{thm:cubic-kuranishi}
For the reconstruction above,
\[
\Lambda\circ\cD_p|_{\mathcal K_B}=0,
\qquad
\Lambda(\zeta)=8.
\]
Consequently,
\[
[\zeta]\ne0
\quad\text{in}\quad
\coker\bigl(\cD_p|_{\mathcal K_B}\bigr).
\]
For a transferred minimal model of the presentation DGLA, the homogeneous Kuranishi terms on
\[
H^1=\C[\alpha]\oplus\C[\xi]
\]
may therefore be normalized as
\[
\boxed{
\kappa_2=0,
\qquad
\kappa_3(u[\alpha]+v[\xi])=v^3[\zeta].
}
\]
In particular,
\[
\kappa_2([\xi],[\xi])=0,
\qquad
\kappa_3([\xi],[\xi],[\xi])\ne0.
\]
\end{theorem}

\begin{proof}
Equation~\eqref{eq:cubic-second-lift} proves that the binary bracket of every pair of tangent classes is zero in cohomology.  For the cubic term, direct matrix substitution into the relation differential gives the two identities involving $\Lambda$; hence $\zeta$ cannot be a strict tangent boundary.

It remains to check the mixed cubic coefficient.  Define $\theta=(S,T)$ by
\begin{align*}
S(\epsilon)&=-2I+N,
&S(\epsilon^2)&=-4N,
&S(\epsilon^3)&=-6N^2,\\
T(\epsilon)&=2I-N,
&T(\epsilon^2)&=4N,
&T(\epsilon^3)&=6N^2.
\end{align*}
Then $\theta\in\mathcal K_B$ and
\[
\cD_p\theta=-\cP_p(\alpha,\eta).
\]
For the tangent vector $u\alpha+v\xi$, the second-order correction may therefore be chosen as $v^2\eta$, and its cubic source is
\[
uv^2\cP_p(\alpha,\eta)+v^3\cP_p(\xi,\eta).
\]
The first summand is a boundary and the second represents $v^3[\zeta]$.  This is precisely the homogeneous cubic part of the minimal Kuranishi map.

Finally, the nonzero class does not depend on the chosen second-order lift.  Any other lift differs from $\eta$ by $s\alpha+t\xi$, and
\[
\cP_p(\xi,\alpha)=0,
\qquad
\cP_p(\xi,\xi)=2\cQ_p(\xi)\in\operatorname{im}\cD_p.
\]
Thus the resulting cubic class is unchanged.
\end{proof}

\begin{corollary}[Second-order extension with third-order failure]\label{cor:cubic-failure}
The tangent direction $\xi$ extends to a strict reconstruction over
\[
\C[t]/(t^3)
\]
by
\[
p_t=p+t\xi+t^2\eta,
\]
but it admits no extension to $\C[t]/(t^4)$.
\end{corollary}

\begin{proof}
The order-two equation is \eqref{eq:cubic-second-lift}.  Any order-three correction would require $\zeta$ to lie in the image of $\cD_p|_{\mathcal K_B}$, contradicting Theorem~\ref{thm:cubic-kuranishi}.
\end{proof}


\subsection{Complete formal germ and the higher Kuranishi tail}\label{subsec:cubic-complete-germ}

The cubic obstruction can in fact be completed exactly.  Choose the exact nine weak-fiber coordinates
\[
w_0,\ldots,w_8
\]
coming from a nullspace basis of the weak observability matrix, normalized so that
\[
\begin{aligned}
\alpha&=(1,0,0,0,0,0,0,0,0),\\
\xi&=(0,1,0,2,0,0,0,0,0).
\end{aligned}
\]
in these coordinates.  Let $u,v$ be tangent coordinates.  Let
$a,b,c,d,e,f,g$ denote transverse coordinates, with
\[
\begin{aligned}
w_0&=u, & w_1&=v+a, & w_2&=b, & w_3&=2v,\\
w_4&=c, & w_5&=d, & w_6&=e, & w_7&=f, & w_8&=g.
\end{aligned}
\]
Thus the linear $u$- and $v$-axes are precisely the tangent classes $[\alpha]$ and $[\xi]$.

\begin{theorem}[Exact local equations of the cubic example]\label{thm:cubic-complete-germ}
After restricting the full strict reconstruction equations to the weak fiber and applying the preceding linear coordinate change, an exact lexicographic Groebner basis is
\[
\boxed{
\begin{gathered}
a+\frac e6,
\qquad
b-\frac e3,
\qquad
c-\frac{2e}{3},
\qquad
d,
f,
g,\\
e^2,
\qquad
ev,
\qquad
(1+2u)e+3v^2,
\qquad
v^3.
\end{gathered}}
\]
Consequently the local and completed local rings at the base reconstruction satisfy
\[
\cO_{\mathbf R_B,p}
\cong
\left(\frac{\C[u,v]}{(v^3)}\right)_{(u,v)},
\qquad
\boxed{
\widehat{\cO}_{\mathbf R_B,p}
\cong
\frac{\C[[u,v]]}{(v^3)}.
}
\]
In particular,
\[
I_\kappa=(v^3),
\qquad
\dim\Obs^{\mathrm{eff}}_p=1.
\]
\end{theorem}

\begin{proof}
The exact weak equations have a nine-dimensional nullspace.  Substitute its basis parameterization into every left, right, and mixed strict relation.  Exact rational Groebner reduction gives the displayed basis; this is verified independently by the accompanying script.

At the origin
\[
q:=1+2u
\]
is a unit.  The relation $qe+3v^2=0$ therefore eliminates the transverse variable:
\[
e=-\frac{3v^2}{q}.
\]
Then $ev=0$ becomes $v^3=0$, while $e^2=0$ becomes a unit multiple of $v^4=0$ and is redundant modulo $v^3$.  The remaining transverse variables are eliminated linearly.  Hence the localized ring is the localization of $\C[u,v]/(v^3)$ at the origin, and completion gives the stated formal ring.  The final assertion follows from the minimal-equation theorem, since $(v^3)/(u,v)(v^3)$ is one-dimensional.
\end{proof}

The theorem already proves that no new independent quartic or higher equation survives.  One can go further and calculate an explicit all-orders Kuranishi representative tied to the cubic certificate \eqref{eq:cubic-certificate}.

\begin{proposition}[Certificate-normalized higher Kuranishi terms]\label{prop:cubic-higher-tail}
Put $q=1+2u$.  Solve the seven transverse equations by
\[
e=-\frac{3v^2}{q},
\qquad
a=-\frac e6,
\qquad b=\frac e3,
\qquad c=\frac{2e}{3},
\qquad d=f=g=0.
\]
Evaluating the full strict residual on this exact formal section and applying $\Lambda/8$ gives the scalar Kuranishi representative
\[
\boxed{
 k_\Lambda(u,v)
 =
 v^3\frac{v+4q}{4q^2}.
}
\]
Its homogeneous pieces begin
\begin{align*}
k_3&=v^3,\\
k_4&=-2uv^3+\frac14v^4,\\
k_5&=4u^2v^3-uv^4,\\
k_6&=-8u^3v^3+3u^2v^4.
\end{align*}
More generally, for $r\ge1$, the part of total degree $3+r$ is
\[
\boxed{
 k_{3+r}
 =
 (-2)^r u^r v^3
 +\frac{r(-2)^{r-1}}{4}u^{r-1}v^4.
}
\]
Every one of these terms is a multiple of the cubic generator.  Since
\[
\frac{4q^2}{v+4q}
\]
is a unit with value $1$ at the origin,
\[
\frac{4q^2}{v+4q}\,k_\Lambda(u,v)=v^3.
\]
Thus an equivalent minimal Kuranishi presentation is exactly
\[
\boxed{
\widetilde\kappa(u[\alpha]+v[\xi])=v^3[\zeta],
\qquad
\widetilde\kappa_n=0\quad(n\ge4).
}
\]
\end{proposition}

\begin{proof}
The exact formal section is obtained by solving the Groebner basis of Theorem~\ref{thm:cubic-complete-germ}.  Direct substitution into all strict relations shows that every residual component is divisible by $v^3$ in $\C[[u,v]]$, and at least one quotient is a unit.  Applying the sparse functional $\Lambda$ gives
\[
\frac18\Lambda(\text{residual})
=
\frac{v^3(v+4q)}{4q^2}.
\]
Expanding
\[
\frac1q=\sum_{n\ge0}(-2u)^n,
\qquad
\frac1{q^2}=\sum_{n\ge0}(n+1)(-2u)^n
\]
gives the displayed homogeneous terms and the general formula.  The final unit rescaling does not change the generated Kuranishi ideal and removes the whole higher tail at once.
\end{proof}

\begin{remark}[Intrinsic content versus presentation tail]
The nonzero cubic class is intrinsic, but the individual quartic and higher coefficients displayed above are not new obstruction classes.  They depend on the chosen formal section and relation generator.  Their simultaneous removal by a unit is exactly the distinction between the intrinsic minimal relation ideal $(v^3)$ and a nonminimal all-orders representative of that same ideal.
\end{remark}

\begin{remark}[Geometric shape of the germ]
The reduced local scheme is the smooth $u$-axis.  Unlike the square-zero example of Section~\ref{sec:two-parameter-example}, where the false direction is pinned to the origin by the mixed equation $uv=0$, the present germ is a uniform cubic thickening of the whole smooth branch:
\[
\Spf\C[[u,v]]/(v^3)
\cong
\Spf\C[[u]]\,\widehat\times\,\Spec\C[v]/(v^3).
\]
The direction $[\xi]$ is therefore a nilpotent normal direction of order three along the branch, not an isolated embedded tangent direction.
\end{remark}

\begin{remark}[Formal normal form versus strict transfer]
The preceding proposition concerns equivalence of minimal Kuranishi presentations.  Such an equivalence may multiply the obstruction generator by a formal unit depending on the tangent variables.  A change of contraction of a fixed DGLA is more restrictive.  The next result shows that, in the present cubic example, the formal pure-cubic normal form cannot be obtained by a strict homotopy transfer from any contraction.
\end{remark}

\subsection{A quartic no-go theorem for strict transferred normal forms}\label{subsec:strict-transfer-no-go}

Retain the tangent basis $\alpha,\xi$ and the weak perturbation space $\mathcal K_B$.  Define the linear map
\[
\mu_\alpha:\mathcal K_B\longrightarrow
\coker\bigl(\cD_p|_{\mathcal K_B}\bigr),
\qquad
k\longmapsto [\cP_p(\alpha,k)].
\]
The freedom to change the $v^3$ coefficient of the transferred $L_\infty$ quasi-isomorphism changes the quartic $uv^3$ source precisely by an element of $\operatorname{im}\mu_\alpha$.

Choose $\eta$ and $\theta$ as above, so that
\[
\cD_p\eta=-\cQ_p(\xi),
\qquad
\cD_p\theta=-\cP_p(\alpha,\eta).
\]
Define the \emph{quartic contraction-normal-form class}
\[
\boxed{
\nu_4(\alpha,\xi)
:=
[\cP_p(\xi,\theta)]
\quad\text{in}\quad
\frac{
\coker(\cD_p|_{\mathcal K_B})
}{\operatorname{im}\mu_\alpha}.
}
\]

\begin{lemma}[Choice independence of $\nu_4$]\label{lem:quartic-choice-independence}
The class $\nu_4(\alpha,\xi)$ is independent of the choices of the second-order lift $\eta$ and the mixed third-order lift $\theta$ satisfying the displayed equations.
\end{lemma}

\begin{proof}
Any other second-order lift is
\[
\eta'=\eta+s\alpha+t\xi.
\]
Because
\[
\cQ_p(\alpha)=0,
\qquad
\cP_p(\alpha,\xi)=0,
\]
one has $\cP_p(\alpha,\eta')=\cP_p(\alpha,\eta)$ exactly.  Hence any corresponding $\theta'$ differs from $\theta$ by a tangent vector $a\alpha+b\xi$.  Therefore
\[
\cP_p(\xi,\theta')-\cP_p(\xi,\theta)
=
 a\cP_p(\xi,\alpha)+b\cP_p(\xi,\xi)
=
2b\cQ_p(\xi),
\]
which lies in $\operatorname{im}\cD_p$ by the defining equation for $\eta$.  Thus the class is unchanged even before quotienting by $\operatorname{im}\mu_\alpha$.
\end{proof}

For a relation vector $\Theta$, define a second sparse functional
\begin{align}\label{eq:quartic-certificate}
\Gamma(\Theta)
={}&
\Theta^L(\epsilon,\epsilon)_{12}
-\Theta^L(\epsilon,\epsilon)_{22}
+\Theta^L(\epsilon,\epsilon)_{23}
-2\Theta^L(\epsilon^2,\epsilon)_{13}\notag\\
&-\Theta^R(\epsilon,\epsilon)_{12}
-\Theta^R(\epsilon,\epsilon)_{23}.
\end{align}

\begin{theorem}[Strict transferred pure-cubic normal form is impossible]\label{thm:strict-transfer-no-go}
For the cubic reconstruction of Section~\ref{sec:cubic-example},
\[
\Gamma\circ\cD_p|_{\mathcal K_B}=0,
\qquad
\Gamma\bigl(\cP_p(\alpha,k)\bigr)=0
\quad(k\in\mathcal K_B),
\]
while
\[
\boxed{
\Gamma\bigl(\cP_p(\xi,\theta)\bigr)=-16.
}
\]
Consequently
\[
\boxed{
\nu_4(\alpha,\xi)\ne0.
}
\]
In particular, for every contraction of the presentation DGLA onto cohomology, the transferred quartic operation is nonzero on the $H^1$-generated sector.  Equivalently,
\[
\boxed{
\ell_4\big|_{\operatorname{Sym}^4 H^1}\ne0
}
\]
for every strictly transferred minimal structure arising from a contraction.  Hence there is no contraction for which
\[
\ell_n\big|_{\operatorname{Sym}^nH^1}=0
\qquad(n\ge4).
\]
\end{theorem}

\begin{proof}
Let $(i,\pi,h)$ be any contraction.  Since the presentation DGLA has no degree-zero term, $H^1$ identifies with $\ker\cD_p$, so the linear term of the transferred Maurer--Cartan section is
\[
x=u\alpha+v\xi.
\]
The quadratic source is $v^2\cQ_p(\xi)$.  Therefore its quadratic correction has the form
\[
F_2=v^2E,
\qquad
E=\eta+s\alpha+t\xi
\]
for some scalars $s,t$.  The coefficient of $uv^2$ in the cubic source is
\[
\cP_p(\alpha,E)=\cP_p(\alpha,\eta),
\]
so the corresponding cubic correction has coefficient
\[
T=\theta+a\alpha+b\xi
\]
for some $a,b$.  Write $K\in\mathcal K_B$ for the coefficient of $v^3$ in the cubic correction; its precise value depends on the chosen representative of the cubic obstruction class and hence on the contraction.

The coefficient of $uv^3$ in the quartic Maurer--Cartan source is then
\[
\Omega_{4,h}
=
\cP_p(\xi,T)+\cP_p(\alpha,K).
\]
The tangent ambiguity in $T$ changes the first summand only by
\[
a\cP_p(\xi,\alpha)+b\cP_p(\xi,\xi)
=2b\cQ_p(\xi)\in\operatorname{im}\cD_p.
\]
Thus the image of $[\Omega_{4,h}]$ in the quotient by $\operatorname{im}\mu_\alpha$ is exactly $\nu_4(\alpha,\xi)$, independently of the contraction.

Exact rational substitution gives the three identities involving $\Gamma$.  Hence $\Gamma$ descends to the dual of
\[
\coker(\cD_p|_{\mathcal K_B})/\operatorname{im}\mu_\alpha
\]
and detects $\nu_4$ with value $-16$.  Therefore $\Omega_{4,h}$ is never a boundary.  The $uv^3$ coefficient of the quartic Kuranishi term is consequently nonzero for every contraction; up to the standard nonzero factorial normalization this is the component $\ell_4(\alpha,\xi,\xi,\xi)$.  Hence the transferred quartic operation cannot vanish identically.
\end{proof}

\begin{proposition}[Exact rank check for the contraction ambiguity]\label{prop:quartic-rank-check}
The image of $\mu_\alpha$ has dimension three.  Moreover,
\[
\rank\bigl(\cD_p,\mu_\alpha,\cP_p(\xi,\theta)\bigr)
=
\rank\bigl(\cD_p,\mu_\alpha\bigr)+1.
\]
Equivalently, the finite linear system obtained by allowing every weak $v^3$ cubic correction and every boundary lift has coefficient rank $20$ and augmented rank $21$.
\end{proposition}

\begin{proof}
These are exact rational rank computations in the nine-dimensional weak fiber.  The sparse functional $\Gamma$ is an independent certificate for the rank jump: it annihilates the first two summands and is nonzero on the last vector.
\end{proof}

\begin{remark}[Why the formal normal form still exists]
There is no contradiction with Proposition~\ref{prop:cubic-higher-tail}.  The formal rescaling
\[
\frac{4(1+2u)^2}{v+4(1+2u)}
\]
acts on the \emph{relation generator} by a tangent-dependent unit.  Such a nonlinear change is permitted when passing between equivalent minimal Kuranishi presentations, but it is not generated by changing only the linear contraction data of the fixed presentation DGLA.  The nonzero class $\nu_4$ measures precisely this gap in the present example.
\end{remark}

\begin{remark}[What remains open after the presentation calculation]
The strict transferred normal-form question is settled negatively for the \emph{presentation} DGLA of this cubic example.  It does not yet follow that the same quartic class is intrinsic to the full constrained Hochschild DGLA, because the latter has a nontrivial differential out of degree two.  The next subsection resolves this distinction explicitly.
\end{remark}

\subsection{The presentation quartic class does not lift intrinsically}\label{subsec:intrinsic-quartic}

Write
\[
\mathsf a,\mathsf x\in\ghoch^1_{B,p}
\]
for the normalized Hochschild one-cocycles corresponding to the tangent vectors
$\alpha,\xi$.  Thus, for $a,b\in A$,
\[
\mathsf a(a\otimes b^{\mathrm{op}})
=
X_\alpha(a)R(b)+L(a)Y_\alpha(b),
\]
and similarly for $\mathsf x$.
Let $\eta=(U,W)$ be the second-order action correction of
\eqref{eq:cubic-second-lift}, and define the normalized Hochschild one-cochain
\[
\mathsf e(a\otimes b^{\mathrm{op}})
=
U(a)R(b)+L(a)W(b)+X_\xi(a)Y_\xi(b).
\]
It is the coefficient of $v^2$ in the enveloping-algebra representation
associated with the second-order action deformation.

Exact cochain calculation gives the stronger identities
\begin{equation}\label{eq:intrinsic-low-identities}
\mathsf a\smile\mathsf a=0,
\qquad
[\mathsf a,\mathsf x]=0,
\qquad
 d_\mu\mathsf e=-\mathsf x\smile\mathsf x,
\qquad
[\mathsf a,\mathsf e]=-2\,\mathsf x\smile\mathsf x.
\end{equation}
In particular, the vanishing of the $u^2$ and $uv$ quadratic sources occurs already at the intrinsic cochain level; a contraction cannot insert arbitrary tangent-valued corrections into those zero sources.

Put
\[
\mathsf s=[\mathsf x,\mathsf e]\in C^2(S,E_\mu).
\]
Then $d_\mu\mathsf s=0$.  We next choose a specific constrained one-cochain $\mathsf k$.
Let $E_{ij}$ denote the standard matrix unit in $M_3(\C)$.  On the two factor axes set
\begin{align*}
X_{\mathsf k}(\epsilon)&=-\tfrac12E_{32},
&X_{\mathsf k}(\epsilon^2)&=-E_{33},
&X_{\mathsf k}(\epsilon^3)&=0,\\
Y_{\mathsf k}(\epsilon)&=\tfrac12E_{32},
&Y_{\mathsf k}(\epsilon^2)&=E_{33},
&Y_{\mathsf k}(\epsilon^3)&=0,
\end{align*}
and set
\[
\mathsf k(\epsilon^i\otimes(\epsilon^j)^{\mathrm{op}})=0
\qquad(i,j\ge1).
\]
The factor-axis pair is a weak $B$-compatible perturbation, hence
$\mathsf k\in\ghoch^1_{B,p}$.
Define
\[
\widehat\zeta:=\mathsf s+d_\mu\mathsf k.
\]
Then $\widehat\zeta$ is closed and represents the same cubic Hochschild class as
$\mathsf s$.  Moreover,
\begin{equation}\label{eq:intrinsic-zeta-projection}
\Lambda\bigl(\Psi^2(\widehat\zeta)\bigr)=8,
\end{equation}
so $[\widehat\zeta]\ne0$ in $H^2(\ghoch_{B,p})$.

\begin{definition}[Acyclic-intersection criterion]\label{def:acyclic-intersection}
Let $(C^\bullet,d)$ be a cochain complex and $W\subset C^n$ a finite-dimensional subspace of source cochains that one wants a contraction projection to annihilate.  We say that $W$ is \emph{acyclicly placeable} if
\[
W\cap Z^n\subseteq B^n.
\]
Equivalently, the image of $W\cap Z^n$ in $H^n(C)$ is zero.
\end{definition}

\begin{lemma}[Placement in a contraction complement]\label{lem:acyclic-placement}
If $W\subset C^n$ is acyclicly placeable, then one can choose a standard vector-space contraction of $C^\bullet$ onto its cohomology for which the degree-$n$ projection vanishes on $W$.
\end{lemma}

\begin{proof}
Choose splittings
\[
Z^n=B^n\oplus H^n_{\mathrm{rep}}.
\]
Because $W\cap Z^n\subseteq B^n$, the image of $W$ in $C^n/Z^n$ has the same dimension as the image of $W/(W\cap B^n)$.  Hence a complement $C^n_{\mathrm{ac}}$ of $Z^n$ may be chosen to contain a complement of $W\cap B^n$ in $W$.  Then
\[
C^n=B^n\oplus H^n_{\mathrm{rep}}\oplus C^n_{\mathrm{ac}}
\]
has projection onto $H^n_{\mathrm{rep}}$ equal to zero on $W$.  Extending such splittings degreewise gives the usual contraction.
\end{proof}

\begin{theorem}[Intrinsic quartic vanishing contraction]\label{thm:intrinsic-ell4-vanish}
For the cubic reconstruction of Section~\ref{sec:cubic-example}, there exists a contraction of the intrinsic constrained Hochschild DGLA onto cohomology such that
\[
\boxed{
\ell_4\big|_{\operatorname{Sym}^4H^1}=0.
}
\]
The same contraction retains the nonzero cubic obstruction, which may be normalized as
\[
\ell_3(u[\mathsf a]+v[\mathsf x],\ldots)
\sim v^3[\widehat\zeta].
\]
Consequently the nonzero presentation class $\nu_4(\alpha,\xi)$ of
Theorem~\ref{thm:strict-transfer-no-go} does not lift to a contraction-independent quartic normal-form obstruction in the intrinsic constrained Hochschild DGLA.
\end{theorem}

\begin{proof}
First note that the two boundaries
\[
\mathsf x\smile\mathsf x,
\qquad
d_\mu\mathsf k
\]
are linearly independent.  Hence the degree-one splitting may be chosen so that both $\mathsf e$ and $\mathsf k$ lie in its acyclic complement, with
\[
h(\mathsf x\smile\mathsf x)=-\mathsf e.
\]
Choose the cubic cohomology representative to be
\[
i[\widehat\zeta]=\widehat\zeta=\mathsf s+d_\mu\mathsf k.
\]
Then the standard contraction identity gives
\[
h(\mathsf s)=-\mathsf k.
\]
Using \eqref{eq:intrinsic-low-identities}, the transferred Maurer--Cartan section through cubic order can therefore be taken as
\[
F_1=u\mathsf a+v\mathsf x,
\qquad
F_2=v^2\mathsf e,
\qquad
F_3=-2uv^2\mathsf e+v^3\mathsf k.
\]
Its cubic residual is exactly $v^3\widehat\zeta$.

The quartic source has no $u^4$ or $u^3v$ term.  Its remaining coefficients are
\begin{align*}
Y_{22}&=[\mathsf a,-2\mathsf e]
      =4\,\mathsf x\smile\mathsf x,\\
Y_{13}&=[\mathsf a,\mathsf k]-2[\mathsf x,\mathsf e],\\
Y_{04}&=[\mathsf x,\mathsf k]+\mathsf e\smile\mathsf e.
\end{align*}
The first is a boundary.  The other two are not closed.  More strongly, their degree-three differentials are linearly independent.  Evaluating them on
\[
(1\otimes\epsilon^{\mathrm{op}})^{\otimes3}
\]
and taking the $(1,2)$ and $(1,3)$ matrix entries gives the exact $2\times2$ minor
\[
\begin{pmatrix}
0&-\tfrac12\\[2pt]
\tfrac12&\tfrac12
\end{pmatrix},
\qquad
\det=\frac14\ne0.
\]
Therefore
\[
\Span\{Y_{22},Y_{13},Y_{04}\}\cap Z^2
\subseteq B^2.
\]
By Lemma~\ref{lem:acyclic-placement}, the degree-two contraction projection can be chosen to annihilate all three quartic source coefficients.  This gives
$\ell_4|_{\operatorname{Sym}^4H^1}=0$.

Finally, the cubic class remains nonzero by \eqref{eq:intrinsic-zeta-projection}.  Thus the vanishing of the intrinsic quartic transfer is not obtained by killing the cubic obstruction.
\end{proof}

\begin{proposition}[Why the presentation class survives]\label{prop:intrinsic-presentation-gap}
For the contraction above, the quartic Hochschild cochains satisfy
\[
\Gamma\bigl(\Psi^2(Y_{13})\bigr)=-16,
\qquad
\Gamma\bigl(\Psi^2(Y_{04})\bigr)=-\frac72.
\]
Thus the two intrinsic quartic sources can lie in an acyclic Hochschild complement even though their presentation images are detected by the sparse presentation certificate.
\end{proposition}

\begin{proof}
Both identities are exact substitutions.  There is no contradiction with Theorem~\ref{thm:intrinsic-ell4-vanish}: the intrinsic sources $Y_{13},Y_{04}$ are not Hochschild cocycles, whereas the presentation DGLA has no degree-three term.  Since
\[
\Psi^3=0,
\]
the degree-three differentials that make these directions acyclic are forgotten by the finite relation model.  Their images therefore become legitimate degree-two presentation classes even though the source cochains may be placed entirely in the intrinsic acyclic complement.
\end{proof}

\begin{corollary}[Failure of a strict intrinsic lift of $\nu_4$]\label{cor:no-intrinsic-nu4-lift}
The presentation class $\nu_4(\alpha,\xi)$ has no lift whose nonvanishing is a contraction-independent obstruction to
\[
\ell_4|_{\operatorname{Sym}^4H^1}=0
\]
in the constrained Hochschild DGLA.  In particular, the presentation quartic no-go is a truncation-sensitive normal-form phenomenon rather than an intrinsic arity-four obstruction of the full Hochschild deformation complex.
\end{corollary}

\begin{remark}[The correct intrinsic arity-four viewpoint]
For the full Hochschild complex the relevant chain-level datum is not an arbitrary quotient of $C^2$ by presentation relations.  It is the harmonic intersection
\[
\frac{W_4\cap Z^2}{W_4\cap B^2},
\]
where $W_4$ is the span of the quartic source cochains produced by the chosen lower-order transfer.  Vanishing of this quotient is precisely what allows $W_4$ to be placed in the acyclic complement.  In the contraction constructed above this quotient is zero.
\end{remark}


\subsection{The intrinsic contraction extends through arity five}\label{subsec:intrinsic-quintic}

The arity-four theorem leaves genuine freedom in the values of the contracting homotopy on the two nonclosed quartic source cochains.  To continue one \emph{single} contraction, this freedom must be chosen compatibly with both the already selected quartic acyclic directions and the new quintic sources.  We now give an explicit compatible choice.

Let $E_{ij}$ again denote the standard matrix units.  Define constrained one-cochains $\mathsf p,\mathsf q\in\ghoch^1_{B,p}$, with zero mixed values, by
\begin{align*}
X_{\mathsf p}(\epsilon)
&=-\tfrac13E_{21}+\tfrac16E_{22},
&Y_{\mathsf p}(\epsilon)
&=\tfrac13E_{21}-\tfrac16E_{22},\\
X_{\mathsf p}(\epsilon^2)
&=-\tfrac23E_{22},
&Y_{\mathsf p}(\epsilon^2)
&=\tfrac23E_{22},\\
X_{\mathsf p}(\epsilon^3)
&=-E_{23},
&Y_{\mathsf p}(\epsilon^3)
&=E_{23},
\end{align*}
and
\begin{align*}
X_{\mathsf q}(\epsilon)
&=-\tfrac12E_{22},
&Y_{\mathsf q}(\epsilon)
&=\tfrac12E_{22},\\
X_{\mathsf q}(\epsilon^2)
&=-E_{23},
&Y_{\mathsf q}(\epsilon^2)
&=E_{23},\\
X_{\mathsf q}(\epsilon^3)
&=0,
&Y_{\mathsf q}(\epsilon^3)
&=0.
\end{align*}
Exact calculation gives
\begin{equation}\label{eq:arity5-degree1-rank}
\rank\Span\{d\mathsf e,d\mathsf k,d\mathsf p,d\mathsf q\}=4,
\end{equation}
so the degree-one acyclic complement may be chosen to contain all four cochains simultaneously.

Set
\begin{equation}\label{eq:arity5-quartic-complement}
C_{13}=Y_{13}-d\mathsf p,
\qquad
C_{04}=Y_{04}-d\mathsf q.
\end{equation}
Then the quartic source decomposition is compatible with
\[
h(Y_{13})=\mathsf p,
\qquad
h(Y_{04})=\mathsf q,
\]
while $C_{13}$ and $C_{04}$ are placed in the degree-two acyclic complement.  Since $Y_{22}=4\,\mathsf x\smile\mathsf x$ and $h(\mathsf x\smile\mathsf x)=-\mathsf e$, the arity-four Taylor coefficient of the transferred section can be chosen as
\begin{equation}\label{eq:intrinsic-F4}
F_4
=
4u^2v^2\mathsf e
-uv^3\mathsf p
-v^4\mathsf q.
\end{equation}

The quintic Maurer--Cartan source
\[
R_5=[F_1,F_4]+[F_2,F_3]
\]
has only the four monomial coefficients
\begin{align}
Z_{32}&=-8\,\mathsf x\smile\mathsf x,\label{eq:arity5-Z32}\\
Z_{23}&=4\mathsf s-[\mathsf a,\mathsf p],\label{eq:arity5-Z23}\\
Z_{14}&=-[\mathsf a,\mathsf q]-[\mathsf x,\mathsf p]-4\mathsf e\smile\mathsf e,\label{eq:arity5-Z14}\\
Z_{05}&=-[\mathsf x,\mathsf q]+[\mathsf e,\mathsf k].\label{eq:arity5-Z05}
\end{align}
Here $Z_{32}$ is a boundary because
\[
Z_{32}=d(8\mathsf e).
\]
The remaining question is whether the nonclosed quartic complement directions and the three non-boundary-type quintic sources can be put into the \emph{same} acyclic complement.

\begin{theorem}[Intrinsic quintic vanishing extension]\label{thm:intrinsic-ell5-vanish}
For the cubic reconstruction of Section~\ref{sec:cubic-example}, the intrinsic contraction of Theorem~\ref{thm:intrinsic-ell4-vanish} can be chosen and extended so that
\[
\boxed{
\ell_4\big|_{\operatorname{Sym}^4H^1}=0,
\qquad
\ell_5\big|_{\operatorname{Sym}^5H^1}=0.
}
\]
The cubic obstruction $[\widehat\zeta]$ remains nonzero.
\end{theorem}

\begin{proof}
Let
\[
W_{4,5}^{\mathrm{nc}}
=
\Span\{C_{13},C_{04},Z_{23},Z_{14},Z_{05}\}
\subset C^2(S,E_\mu).
\]
We show that the Hochschild differential is injective on this five-dimensional space.

Write
\[
x=\epsilon\otimes1,
\qquad
y=1\otimes\epsilon^{\mathrm{op}},
\qquad
m=\epsilon\otimes\epsilon^{\mathrm{op}}.
\]
For a normalized three-cochain $\Phi$, evaluate successively the coordinates
\[
\Phi(y,y,y)_{12},\quad
\Phi(y,y,y)_{13},\quad
\Phi(y,y,y)_{23},\quad
\Phi(y,y,m)_{13},\quad
\Phi(y,x,y)_{13}.
\]
On the ordered columns
\[
dC_{13},\ dC_{04},\ dZ_{23},\ dZ_{14},\ dZ_{05},
\]
these five functionals give the exact matrix
\begin{equation}\label{eq:arity5-minor}
\begin{pmatrix}
0&-\tfrac12&0&-\tfrac23&-\tfrac12\\
\tfrac12&\tfrac12&-\tfrac13&\tfrac13&-1\\
0&1&0&\tfrac13&0\\
0&0&0&\tfrac13&0\\
\tfrac12&-\tfrac12&\tfrac13&-\tfrac13&0
\end{pmatrix},
\qquad
\det=\frac1{18}\ne0.
\end{equation}
Hence
\[
\rank\Span\{dC_{13},dC_{04},dZ_{23},dZ_{14},dZ_{05}\}=5,
\]
and therefore
\[
W_{4,5}^{\mathrm{nc}}\cap Z^2=0.
\]
The degree-two acyclic complement can consequently be chosen to contain the whole space $W_{4,5}^{\mathrm{nc}}$ at once.  Together with the boundary $Z_{32}$, every quintic source coefficient has zero cohomology projection.  The same splitting already kills all quartic coefficients by construction, so the resulting single contraction satisfies
\[
\ell_4|_{\operatorname{Sym}^4H^1}
=
\ell_5|_{\operatorname{Sym}^5H^1}
=0.
\]
Nothing in this extension changes the cubic representative $\widehat\zeta$, whose nonvanishing was established by \eqref{eq:intrinsic-zeta-projection}.
\end{proof}

\begin{remark}[No arity-five obstruction for this intrinsic splitting]
The arity-five calculation is stronger than checking the quintic sources in isolation.  The determinant \eqref{eq:arity5-minor} tests the \emph{joint} span of the previously fixed quartic acyclic directions and the new quintic directions.  Thus the argument really constructs one compatible contraction through arity five; it does not silently replace the contraction between orders.
\end{remark}


\subsection{The intrinsic contraction extends through arity six}\label{subsec:intrinsic-sextic}

We now continue the \emph{same} contraction one order further.  The arity-five theorem fixed the degree-two complement directions but left freedom in the values of the homotopy on the three non-boundary-type quintic sources.  We choose explicit constrained one-cochains
\[
\mathsf r,\mathsf w,\mathsf t\in\ghoch^1_{B,p}
\]
with the following values.  The cochain $\mathsf r$ has only factor-axis values
\begin{align*}
X_{\mathsf r}(\epsilon)
&=-\tfrac13E_{31}+\tfrac16E_{32},
&Y_{\mathsf r}(\epsilon)
&=\tfrac13E_{31}-\tfrac16E_{32},\\
X_{\mathsf r}(\epsilon^2)
&=-\tfrac23E_{32},
&Y_{\mathsf r}(\epsilon^2)
&=\tfrac23E_{32},\\
X_{\mathsf r}(\epsilon^3)
&=-E_{33},
&Y_{\mathsf r}(\epsilon^3)
&=E_{33}.
\end{align*}
The cochains $\mathsf w$ and $\mathsf t$ have zero factor-axis values and are supported only at
\[
m=\epsilon\otimes\epsilon^{\mathrm{op}},
\]
where
\[
\mathsf w(m)=E_{11},
\qquad
\mathsf t(m)=E_{12}.
\]
Exact calculation gives
\begin{equation}\label{eq:arity6-degree1-rank}
\rank\Span\{d\mathsf e,d\mathsf k,d\mathsf p,d\mathsf q,d\mathsf r,d\mathsf w,d\mathsf t\}=7.
\end{equation}
Thus all seven cochains may be included simultaneously in one degree-one acyclic complement.

Set
\begin{equation}\label{eq:arity6-quintic-complement}
D_{23}=Z_{23}-d\mathsf r,
\qquad
D_{14}=Z_{14}-d\mathsf w,
\qquad
D_{05}=Z_{05}-d\mathsf t.
\end{equation}
This corresponds to choosing
\[
h(Z_{23})=\mathsf r,
\qquad
h(Z_{14})=\mathsf w,
\qquad
h(Z_{05})=\mathsf t.
\]
Since $Z_{32}=d(8\mathsf e)$, the arity-five Taylor coefficient can then be chosen as
\begin{equation}\label{eq:intrinsic-F5}
F_5
=
-8u^3v^2\mathsf e
-u^2v^3\mathsf r
-uv^4\mathsf w
-v^5\mathsf t.
\end{equation}

The sextic Maurer--Cartan source is
\[
R_6=[F_1,F_5]+[F_2,F_4]+\frac12[F_3,F_3].
\]
It has five monomial coefficients:
\begin{align}
U_{42}&=-8[\mathsf a,\mathsf e],\label{eq:arity6-U42}\\
U_{33}&=-[\mathsf a,\mathsf r]-8\mathsf s,\label{eq:arity6-U33}\\
U_{24}&=-[\mathsf a,\mathsf w]-[\mathsf x,\mathsf r]+12\,\mathsf e\smile\mathsf e,\label{eq:arity6-U24}\\
U_{15}&=-[\mathsf a,\mathsf t]-[\mathsf x,\mathsf w]-[\mathsf e,\mathsf p]-2[\mathsf e,\mathsf k],\label{eq:arity6-U15}\\
U_{06}&=-[\mathsf x,\mathsf t]-[\mathsf e,\mathsf q]+\mathsf k\smile\mathsf k.\label{eq:arity6-U06}
\end{align}
Here $\mathsf s=[\mathsf x,\mathsf e]$ is the nonzero cubic source fixed in the intrinsic construction above.  The coefficient $U_{42}$ is a boundary.  Indeed,
\[
[\mathsf a,\mathsf e]=-2\,\mathsf x\smile\mathsf x,
\qquad
d\mathsf e=-\mathsf x\smile\mathsf x,
\]
so
\begin{equation}\label{eq:arity6-U42-boundary}
U_{42}=d(-16\mathsf e).
\end{equation}

\begin{theorem}[Intrinsic sextic vanishing extension]\label{thm:intrinsic-ell6-vanish}
For the cubic reconstruction of Section~\ref{sec:cubic-example}, one intrinsic contraction can be chosen so that
\[
\boxed{
\ell_4\big|_{\operatorname{Sym}^4H^1}
=
\ell_5\big|_{\operatorname{Sym}^5H^1}
=
\ell_6\big|_{\operatorname{Sym}^6H^1}
=0.
}
\]
The cubic obstruction $[\widehat\zeta]$ remains nonzero.
\end{theorem}

\begin{proof}
Let
\[
W_{4,5,6}^{\mathrm{nc}}
=
\Span\{C_{13},C_{04},D_{23},D_{14},D_{05},U_{33},U_{24},U_{15},U_{06}\}.
\]
We prove that the Hochschild differential is injective on this nine-dimensional space.  As before set
\[
x=\epsilon\otimes1,
\qquad
y=1\otimes\epsilon^{\mathrm{op}},
\qquad
m=\epsilon\otimes\epsilon^{\mathrm{op}}.
\]
Evaluate degree-three cochains successively at the nine coordinates
\[
(yyy)_{11},\ (yyy)_{12},\ (yyy)_{13},\ (yyy)_{23},\ (yyy)_{32},\
(yym)_{13},\ (yym)_{23},\ (y\,y^2\,y)_{13},\ (yxy)_{13}.
\]
On the ordered columns
\[
dC_{13},dC_{04},dD_{23},dD_{14},dD_{05},dU_{33},dU_{24},dU_{15},dU_{06},
\]
these functionals give the exact matrix
\begin{equation}\label{eq:arity6-minor}
\begin{pmatrix}
0&0&0&0&0&0&\tfrac13&0&0\\
0&-\tfrac12&0&-\tfrac23&-\tfrac12&-\tfrac13&-\tfrac12&1&0\\
\tfrac12&\tfrac12&-\tfrac13&\tfrac13&-1&-\tfrac16&\tfrac16&0&0\\
0&1&0&\tfrac13&0&-\tfrac13&\tfrac13&-\tfrac13&0\\
0&0&0&0&0&0&0&0&-\tfrac12\\
0&0&0&\tfrac13&0&\tfrac13&0&0&0\\
0&0&0&0&0&0&\tfrac13&-\tfrac13&0\\
0&0&0&-1&-1&-\tfrac13&-\tfrac23&2&1\\
\tfrac12&-\tfrac12&\tfrac13&-\tfrac13&0&\tfrac16&\tfrac56&0&1
\end{pmatrix},
\qquad
\det=-\frac1{486}\ne0.
\end{equation}
Hence
\[
\rank\Span\{dC_{13},dC_{04},dD_{23},dD_{14},dD_{05},dU_{33},dU_{24},dU_{15},dU_{06}\}=9,
\]
and therefore
\[
W_{4,5,6}^{\mathrm{nc}}\cap Z^2=0.
\]
The degree-two acyclic complement can consequently be chosen to contain all nine directions simultaneously.  Together with the boundary coefficient $U_{42}$, every sextic source coefficient has zero cohomology projection.  The same degree-one and degree-two splittings already realize the quartic and quintic vanishings, so this single contraction satisfies the claimed identities through arity six.  None of these choices alters the previously fixed nonzero cubic class $[\widehat\zeta]$.
\end{proof}

\begin{remark}[No intrinsic obstruction through arity six]\label{rem:no-obstruction-through-six}
The arity-six test again uses the joint source space rather than the sextic coefficients in isolation.  Thus Theorem~\ref{thm:intrinsic-ell6-vanish} constructs one compatible contraction through arity six.  It does \emph{not} prove that all higher transferred operations vanish.  Arity seven is the next point at which a new closed non-boundary combination may appear inside the growing joint acyclic source space.
\end{remark}


\subsection{The intrinsic contraction extends through arity seven}\label{subsec:intrinsic-arity7}

We keep every splitting choice from the arity-six construction and choose homotopy lifts for the four nonclosed sextic coefficients.  Let
\[
m=\epsilon\otimes\epsilon^{\mathrm{op}}.
\]
Define four constrained one-cochains $\mathsf A,\mathsf B,\mathsf C,\mathsf D$ with zero factor-axis values, supported only at $m$, by
\[
\mathsf A(m)=E_{21},\qquad
\mathsf B(m)=E_{22},\qquad
\mathsf C(m)=E_{23},\qquad
\mathsf D(m)=E_{31}.
\]
Exact calculation gives
\begin{equation}\label{eq:arity7-degree1-rank}
\rank\Span\{d\mathsf e,d\mathsf k,d\mathsf p,d\mathsf q,d\mathsf r,d\mathsf w,d\mathsf t,
 d\mathsf A,d\mathsf B,d\mathsf C,d\mathsf D\}=11.
\end{equation}
Thus all eleven degree-one homotopy directions can be included simultaneously in one acyclic splitting.

Set
\begin{equation}\label{eq:arity7-sextic-complement}
G_{33}=U_{33}-d\mathsf A,\qquad
G_{24}=U_{24}-d\mathsf B,\qquad
G_{15}=U_{15}-d\mathsf C,\qquad
G_{06}=U_{06}-d\mathsf D.
\end{equation}
Accordingly the arity-six Taylor coefficient can be chosen as
\begin{equation}\label{eq:intrinsic-F6}
F_6
=
16u^4v^2\mathsf e
-u^3v^3\mathsf A
-u^2v^4\mathsf B
-uv^5\mathsf C
-v^6\mathsf D.
\end{equation}

The arity-seven Maurer--Cartan source is
\[
R_7=[F_1,F_6]+[F_2,F_5]+[F_3,F_4].
\]
Its six monomial coefficients are
\begin{align}
V_{52}&=16[\mathsf a,\mathsf e],\label{eq:arity7-V52}\\
V_{43}&=-[\mathsf a,\mathsf A]+16[\mathsf x,\mathsf e],\label{eq:arity7-V43}\\
V_{34}&=-[\mathsf a,\mathsf B]-[\mathsf x,\mathsf A]-16[\mathsf e,\mathsf e],\label{eq:arity7-V34}\\
V_{25}&=-[\mathsf a,\mathsf C]-[\mathsf x,\mathsf B]-[\mathsf e,\mathsf r]
+2[\mathsf e,\mathsf p]+4[\mathsf k,\mathsf e],\label{eq:arity7-V25}\\
V_{16}&=-[\mathsf a,\mathsf D]-[\mathsf x,\mathsf C]-[\mathsf e,\mathsf w]
+2[\mathsf e,\mathsf q]-[\mathsf k,\mathsf p],\label{eq:arity7-V16}\\
V_{07}&=-[\mathsf x,\mathsf D]-[\mathsf e,\mathsf t]-[\mathsf k,\mathsf q].\label{eq:arity7-V07}
\end{align}
The leading coefficient is again a boundary.  Since $[\mathsf a,\mathsf e]=-2\mathsf x\smile\mathsf x$ and $d\mathsf e=-\mathsf x\smile\mathsf x$,
\begin{equation}\label{eq:arity7-V52-boundary}
V_{52}=d(32\mathsf e).
\end{equation}

\begin{theorem}[Intrinsic arity-seven vanishing extension]\label{thm:intrinsic-ell7-vanish}
For the cubic reconstruction of Section~\ref{sec:cubic-example}, one intrinsic contraction can be chosen so that
\[
\boxed{
\ell_4\big|_{\operatorname{Sym}^4H^1}
=
\ell_5\big|_{\operatorname{Sym}^5H^1}
=
\ell_6\big|_{\operatorname{Sym}^6H^1}
=
\ell_7\big|_{\operatorname{Sym}^7H^1}
=0.
}
\]
The cubic obstruction $[\widehat\zeta]$ remains nonzero.
\end{theorem}

\begin{proof}
Use the actual sextic complement vectors from \eqref{eq:arity7-sextic-complement} and define
\[
W_{4,5,6,7}^{\mathrm{nc}}
=
\Span\{C_{13},C_{04},D_{23},D_{14},D_{05},G_{33},G_{24},G_{15},G_{06},
V_{43},V_{34},V_{25},V_{16},V_{07}\}.
\]
Because $d^2=0$, the differentials of $G_{33},G_{24},G_{15},G_{06}$ are the same as those of $U_{33},U_{24},U_{15},U_{06}$ used in the arity-six certificate.  It therefore suffices to prove that the five new degree-seven directions increase the joint differential rank from nine to fourteen.

As before write
\[
x=\epsilon\otimes1,\qquad y=1\otimes\epsilon^{\mathrm{op}},\qquad m=\epsilon\otimes\epsilon^{\mathrm{op}}.
\]
On the ordered columns
\[
dC_{13},dC_{04},dD_{23},dD_{14},dD_{05},dG_{33},dG_{24},dG_{15},dG_{06},
dV_{43},dV_{34},dV_{25},dV_{16},dV_{07},
\]
consider the fourteen coordinate evaluations
\[
\begin{gathered}
(yyy)_{11},(yyy)_{12},(yyy)_{13},(yyy)_{21},(yyy)_{22},(yyy)_{23},(yyy)_{32},\\
(yyx)_{11},(yyx)_{12},(yym)_{13},(yym)_{23},(y\,y^2\,y)_{13},(yxy)_{13},(yxy)_{23}.
\end{gathered}
\]
The resulting exact $14\times14$ minor, denoted $M_7$, satisfies
\begin{equation}\label{eq:arity7-minor}
\det M_7=\frac{11}{34992}\ne0.
\end{equation}
Hence
\[
\rank d\big|_{W_{4,5,6,7}^{\mathrm{nc}}}=14,
\qquad
W_{4,5,6,7}^{\mathrm{nc}}\cap Z^2=0.
\]
Thus the degree-two acyclic complement can be enlarged to contain all fourteen directions simultaneously.  Together with the boundary coefficient $V_{52}$, every arity-seven source coefficient has zero cohomology projection.  Equation~\eqref{eq:arity7-degree1-rank} guarantees that the required eleven degree-one homotopy directions coexist in the same splitting.  Therefore the same intrinsic contraction realizes the quartic, quintic, sextic, and arity-seven vanishings while preserving the previously fixed nonzero cubic class.
\end{proof}

\begin{remark}[No intrinsic obstruction through arity seven]\label{rem:no-obstruction-through-seven}
The rank-fourteen certificate is a simultaneous history test: it includes every nonclosed direction already committed at arities four, five, and six.  Thus the arity-seven vanishing is realized by one contraction rather than by changing the splitting at each order.  The calculation still does not imply an all-orders pure-cubic transfer.  Arity eight is the next possible location for a nonzero harmonic-intersection obstruction.
\end{remark}


\section{Finite feedback beyond arity seven}\label{sec:finite-feedback}

The intrinsic transfer calculation of the preceding section can be reorganized as a finite-feedback problem.  The point of this reorganization is not merely computational.  It separates three structures that behave differently at higher order: the harmonic obstruction sector, the finite controller box in degree one, and the actuator directions by which the controller moves future Maurer--Cartan sources out of the closed sector.

Throughout this section we remain in the cubic example of Section~\ref{sec:cubic-example}.  Write
\[
S=A^e\cong \C[x,y]/(x^4,y^4),
\qquad
E=\End(V)\cong M_3(\C),
\]
and denote by $\alpha,\xi\in H^1$ the tangent generators used above.  The low-order homotopy cochains $\mathsf e,\mathsf k,\mathsf p,\ldots$ are those fixed by the intrinsic contraction.

\subsection{A split pure-mixed reservoir}

Let
\[
M_{\mathrm{mix}}
=
\bigoplus_{1\le i,j\le3}E\,\delta_{ij}
\subset \ghoch^1_{B,p}
\]
be the normalized one-cochains that vanish on both factor axes and are supported on mixed monomials $x^iy^j$.  For such a cochain $f$, the Hochschild differential satisfies
\[
(d f)(x^i,y^j)=-f(x^iy^j).
\]
Consequently the restriction map
\[
\rho:\ghoch^2_{B,p}\longrightarrow M_{\mathrm{mix}},
\qquad
(\rho\Phi)(x^iy^j)=-\Phi(x^i,y^j),
\]
satisfies
\begin{equation}\label{eq:split-mixed}
\boxed{\rho\,d|_{M_{\mathrm{mix}}}=\mathrm{id}.}
\end{equation}
Thus the mixed degree-one directions form an explicit split contractible reservoir.  This proves compatibility of individual mixed lifts, but it does not by itself imply that all recursively generated degree-two sources lie in one acyclic complement.  The additional condition required at arity $n$ is
\[
W_{\le n}^{\mathrm{nc}}\cap Z^2=0,
\]
where $W_{\le n}^{\mathrm{nc}}$ denotes the span of the nonclosed source history through that arity.

\subsection{Finite controller boxes and the first bifurcations}

The arity-eight-through-twelve computation enlarges the same intrinsic contraction without producing a new harmonic class.  Exact source-history ranks grow while the nonclosed source directions remain compatible with one acyclic complement.  At arity thirteen, however, a fixed controller box first loses reachability of one harmonic coordinate.

Let $U_{10}$ denote the finite degree-one controller box fixed through arity ten.  The arity-thirteen feedback system has one exact uncontrollable equation,
\begin{equation}\label{eq:arity13-danger}
\boxed{p_\eta(R_{13}^{(7,6)})=5824,}
\end{equation}
while every $U_{10}$-control coefficient in that equation vanishes.  A fresh pure-mixed direction
\[
\mathsf u_*(\epsilon,\epsilon^3)=E_{32}
\]
adds two new acyclic differential directions and rescues the transfer.  Set
\[
U_{11}=U_{10}\oplus\C\mathsf u_*.
\]
The same phenomenon repeats at the next orders, but the required actuator directions evolve.

\begin{table}[ht]
\centering
\small
\begin{tabular}{@{}cccc@{}}
\toprule
arity & finite box before rescue & distinguished $\eta$ source & fresh rescue direction \\
\midrule
13 & $U_{10}$ & $p_\eta(R_{13}^{(7,6)})=5824$ & $(\epsilon,\epsilon^3;E_{32})$ \\
14 & $U_{11}$ & $p_\eta(R_{14}^{(8,6)})=28672$ & $(\epsilon,\epsilon^3;E_{33})$ \\
15 & $U_{12}$ & $p_\eta(R_{15}^{(9,6)})=64000$ & $(\epsilon^2,\epsilon;E_{11})$ \\
16 & $U_{13}$ & $p_\eta(R_{16}^{(10,6)})=-135168$ & $(\epsilon^2,\epsilon;E_{21})$ \\
17 & $U_{14}$ & $p_\eta(R_{17}^{(11,6)})=494592$ & $(\epsilon^2,\epsilon;E_{31})$ \\
18 & $U_{15}$ & $p_\eta(R_{18}^{(12,6)})=-1103872$ & $(\epsilon^2,\epsilon;E_{32})$ \\
\bottomrule
\end{tabular}
\caption{Exact finite-box reachability failures and the fresh directions used to rescue them.  The signs are those of the fixed intrinsic contraction and harmonic projection.}
\label{tab:finite-box-bifurcation}
\end{table}

At every stage in Table~\ref{tab:finite-box-bifurcation}, the enlarged quadratic box remains homologically healthy: its closed non-boundary quotient is still two-dimensional and is detected by the same cubic coordinate $\zeta$ and danger coordinate $\eta$.  Thus the bifurcations are controller-reachability phenomena rather than births of new $H^2$ classes.

\section{A block actuator flag}\label{sec:block-flag}

The cross-support jump at arity fifteen suggests studying the complete mixed matrix block at the support $(\epsilon^2,\epsilon)$.  After the directions $E_{11}$ and $E_{21}$ have entered the controller, let
\[
\mathcal M_{21}^{\mathrm{rem}}
=
\Span\{E_{12},E_{13},E_{22},E_{23},E_{31},E_{32},E_{33}\},
\]
where each matrix unit denotes the pure-mixed one-cochain supported at $(\epsilon^2,\epsilon)$.  Modulo the exact source history $\mathcal H_{16}=dW_{\le16}$, define
\[
\overline T_\alpha(u)=[d[\alpha,u]],
\qquad
\overline T_\xi(u)=[d[\xi,u]].
\]

\begin{theorem}[Block actuator flag]\label{thm:block-actuator-flag}
For the cubic example and the arity-sixteen quotient history,
\[
\ker\overline T_\alpha
=
\Span\{E_{12},E_{13},E_{22},E_{23}\},
\]
\[
\ker\overline T_\xi
=
\C E_{13}.
\]
Moreover
\[
\rank\overline T_\alpha=3,
\qquad
\rank\overline T_\xi=6,
\qquad
\dim(\operatorname{im}\overline T_\alpha+
     \operatorname{im}\overline T_\xi)=8,
\]
and
\begin{equation}\label{eq:block-intersection}
\boxed{
T_\alpha(E_{33})=T_\xi(E_{23})
}
\end{equation}
is an exact cochain identity.  Hence the two image spaces meet in exactly one nonzero line.

Let
\[
I=\C E_{13},
\qquad
K_\alpha=\Span\{E_{12},E_{13},E_{22},E_{23}\}.
\]
For every nonzero $u\in\mathcal M_{21}^{\mathrm{rem}}$, the actuator type is determined by the flag
\[
0\subset I\subset K_\alpha\subset\mathcal M_{21}^{\mathrm{rem}}:
\]
\[
u\in I\Longrightarrow(0,0),
\qquad
u\in K_\alpha\setminus I\Longrightarrow(0,1),
\qquad
u\notin K_\alpha\Longrightarrow(1,1),
\]
and every vector in the last stratum has joint rank two.
\end{theorem}

For the coordinate matrix units the theorem reads
\[
\begin{array}{c|ccccccc}
& E_{12}&E_{13}&E_{22}&E_{23}&E_{31}&E_{32}&E_{33}\\
\hline
\text{type}&(0,1)&(0,0)&(0,1)&(0,1)&(1,1)&(1,1)&(1,1).
\end{array}
\]
This theorem was used predictively rather than only retrospectively.  Before computing arity seventeen it identifies $E_{31}$ as the first unused coordinate direction in the genuinely two-sided stratum.  The arity-seventeen calculation indeed selects $E_{31}$.  After that direction is absorbed into the history, the flag itself remains unchanged and the two-sided coordinate stratum becomes $\{E_{32},E_{33}\}$.  The next calculation selects $E_{32}$.  After the arity-eighteen history update, the stratum becomes $\{E_{33}\}$ and the same kernel flag persists.

Thus, for the two history transitions actually computed,
\begin{equation}\label{eq:deletion-dynamics}
\boxed{
\{E_{31},E_{32},E_{33}\}
\longrightarrow
\{E_{32},E_{33}\}
\longrightarrow
\{E_{33}\}.
}
\end{equation}
The identity~\eqref{eq:block-intersection} survives each update.  After the second deletion the one-dimensional intersection line becomes the entire $\alpha$-image.

\section{Stationary reuse at arities nineteen through twenty-two}\label{sec:stationary-reuse}

The deletion picture does not continue by consuming $E_{33}$ immediately.  At arity nineteen the already-used direction $E_{32}$ has entered the finite box and directly controls the formerly dangerous $v^6$ row.  This begins a stationary-reuse phase.

Let
\[
U_{16}=U_{15}\oplus\C\{E_{32}\text{ at }(\epsilon^2,\epsilon)\}.
\]
This box has dimension $35$ and the quadratic closed non-boundary quotient remains two-dimensional.  The exact feedback matrices at arities nineteen through twenty-two are all of full row rank:
\[
\rank M_{19}^{U_{16}}=28,
\qquad
\rank M_{20}^{U_{16}}=30,
\qquad
\rank M_{21}^{U_{16}}=32,
\qquad
\rank M_{22}^{U_{16}}=34,
\]
with the same ranks after adjoining the right-hand sides.  Thus no finite-box enlargement is required at these four orders.

At each order there is a distinguished $\eta$ equation at fixed $v$-degree six.  Its entire control row has exactly one nonzero actuator coefficient, namely the $E_{32}$ component of the final independent channel.  The actuator sensitivity is stationary:
\begin{equation}\label{eq:stationary-sensitivity}
\boxed{A=1103872.}
\end{equation}
Writing the distinguished equation as
\[
B_n+A c_n=0,
\]
the exact values are
\begin{align*}
B_{19}&=1003520,& c_{19}&=-\frac{10}{11},\\
B_{20}&=-6766592,& c_{20}&=\frac{472}{77},\\
B_{21}&=-13877248,& c_{21}&=\frac{88}{7},\\
B_{22}&=-35192832,& c_{22}&=\frac{17184}{539}.
\end{align*}
A minimal controller uses only $\Span\{\mathsf e,\mathsf k,\mathsf p\}$ on all free channels together with one $E_{32}$ coordinate on the distinguished channel.  Exact pivot-minor certificates have determinants
\[
-9042919424,
\qquad
-18085838848,
\qquad
-36171677696
\]
at arities nineteen, twenty, and twenty-one respectively; at arity twenty-two the corresponding exact pivot determinant is $-72343355392$.

The corresponding source histories have exact ranks
\[
\begin{aligned}
\rank_\Q dW_{\le19}&=126,
&\qquad
\rank_\Q dW_{\le20}&=141,\\
\rank_\Q dW_{\le21}&=157,
&\qquad
\rank_\Q dW_{\le22}&=174.
\end{aligned}
\]
At each of these orders exactly three high-rail source monomials recycle into the previous provenance space.  The leading closed terms alternate sign and remain boundaries:
\[
65536Q_{\alpha\mathsf e}=d(131072\mathsf e),
\]
\[
-131072Q_{\alpha\mathsf e}=d(-262144\mathsf e),
\]
\[
262144Q_{\alpha\mathsf e}=d(524288\mathsf e),
\]
\[
-524288Q_{\alpha\mathsf e}=d(-1048576\mathsf e).
\]
Meanwhile the unused $E_{33}$ direction remains genuinely two-sided modulo every updated history through arity twenty-two.

\section{Controller terminality and the low-rail transfer state}\label{sec:low-rail}

The stationary reuse coefficients suggest a recurrence problem.  A direct filtration argument shows, however, that the $E/K/P$ controller variables are not the memory state that drives future forcing.

\subsection{Controller terminality}

For a polynomial monomial $u^av^b$, write $\nu(u^av^b)=b$.  Every independent stationary feedback channel has $\nu\ge6$.  The bracket adds polynomial exponents.  Since
\[
F_1=u\alpha+v\xi
\]
has $v$-degrees zero and one, while every $F_m$, $m\ge2$, has minimum $v$-degree two on the fixed section, a controller perturbation of $v$-degree at least six cannot enter the forced rails $v^2,v^3,v^4,v^5$.  After its current arity it can pair only with terms of $v$-degree at least two, so every later contribution has $v$-degree at least eight and cannot return to the distinguished $v^6$ row.

\begin{proposition}[Controller terminality]\label{prop:controller-terminality}
Let $\mathcal C_n^{EKP}$ be the stationary $E/K/P$ controller sector at arity $n$, and let
\[
\mathcal O_n:\mathcal C_n^{EKP}\to(B_{n+1},B_{n+2},\ldots)
\]
be the future distinguished-forcing observability map.  On the stationary phase computed here,
\[
\boxed{\mathcal O_n=0.}
\]
Hence the observable-minimal $E/K/P$ state is zero-dimensional.
\end{proposition}

Thus the reuse coefficient $c_n$ is a terminal current-arity correction, not a one-step memory variable.  This explains structurally why
\[
\frac{\partial B_{20}}{\partial c_{19}}=0,
\qquad
\frac{\partial B_{21}}{\partial c_{20}}=0.
\]

\subsection{Autonomous low rails}

The actual memory lies below $v$-degree six.  Define
\begin{equation}\label{eq:low-rails}
\boxed{x_b(k)=F_{k+b}^{(k,b)},\qquad b=2,3,4,5.}
\end{equation}
These cochain-valued sequences satisfy exact triangular quadratic recursions.  With $h$ denoting the fixed intrinsic homotopy,
\begin{equation}\label{eq:x2-rec}
x_2(k)=-h[\alpha,x_2(k-1)],
\qquad
x_2(k)=(-2)^k\mathsf e.
\end{equation}
Next,
\begin{equation}\label{eq:x3-rec}
x_3(k)=-h\bigl([\alpha,x_3(k-1)]+[\xi,x_2(k)]\bigr).
\end{equation}
The $v^4$ rail obeys
\begin{equation}\label{eq:x4-rec}
\begin{aligned}
x_4(k)=-h\Bigg(&[\alpha,x_4(k-1)]+[\xi,x_3(k)]\\
&+\frac12\sum_{a+b=k}[x_2(a),x_2(b)]\Bigg),
\end{aligned}
\end{equation}
and the $v^5$ rail obeys
\begin{equation}\label{eq:x5-rec}
\begin{aligned}
x_5(k)=-h\Bigg(&[\alpha,x_5(k-1)]+[\xi,x_4(k)]\\
&+\frac12\sum_{a+b=k}
\bigl([x_2(a),x_3(b)]+[x_3(b),x_2(a)]\bigr)\Bigg).
\end{aligned}
\end{equation}
All available transitions through arity twenty-one have been verified exactly.

The cochain-mode ranks on the exact trajectory are
\begin{equation}\label{eq:rail-ranks}
\boxed{
\dim W_2=1,
\qquad
\dim W_3=7,
\qquad
\dim W_4=7,
\qquad
\dim W_5=13,
}
\end{equation}
where $W_b=\Span\{x_b(k)\}$ over the computed range.

\subsection{Frozen transfer law and two out-of-sample tests}

Define
\[
q(\lambda)=
\lambda^5-4\lambda^4+12\lambda^3-32\lambda^2+80\lambda-192.
\]
On the initial exact trajectory the four low rails admit the annihilators
\[
p_2(\lambda)=\lambda+2,
\qquad
p_3(\lambda)=(\lambda+2)^2q(\lambda),
\]
\[
p_4(\lambda)=(\lambda+2)^3q(\lambda),
\qquad
p_5(\lambda)=(\lambda+2)^4q(\lambda)^2.
\]
For
\[
Y_k=(x_2(k),x_3(k),x_4(k),x_5(k)),
\]
the first constant scalar recurrence on the fitting window has order fourteen:
\begin{equation}\label{eq:Y14-rec}
\boxed{
\begin{aligned}
Y_{k+14}={}&-589824Y_k-688128Y_{k+1}-200704Y_{k+2}\\
&+1536Y_{k+7}+896Y_{k+8}.
\end{aligned}}
\end{equation}
Its characteristic polynomial is
\begin{equation}\label{eq:low-poly}
\boxed{
p(\lambda)=(\lambda+2)^4q(\lambda)^2.
}
\end{equation}
Equivalently,
\[
p(\lambda)=
\lambda^{14}-896\lambda^8-1536\lambda^7
+200704\lambda^2+688128\lambda+589824.
\]

Let $M_{\mathrm{low}}$ be the $14\times14$ companion matrix of~\eqref{eq:Y14-rec}.  The matrix was frozen before computing further arities.  It predicted $Y_{17}$ and $Y_{18}$ exactly.  The first prediction was then checked by an independent arity-twenty-two quadratic recursion, and the second by an independent arity-twenty-three low-rail recursion.  In both cases every cochain coordinate agreed.  Thus the finite realization passed two genuine out-of-sample tests without refitting.

\subsection{The distinguished forcing as a low-rail observable}

For $n=k+6$, before the current $v^6$ controller is applied, the distinguished forcing is determined entirely by the low rails:
\begin{equation}\label{eq:B-low-observable}
\boxed{
\begin{aligned}
B_{k+6}=p_\eta\Bigg(&[\xi,x_5(k)]\\
&+\frac12\sum_{a+b=k}
\Big([x_2(a),x_4(b)]+[x_4(b),x_2(a)]+[x_3(a),x_3(b)]\Big)
\Bigg).
\end{aligned}}
\end{equation}
For example,
\[
B_{19}=1003520,
\qquad
B_{20}=-6766592,
\qquad
B_{21}=-13877248,
\]
and the stationary $E_{32}$ coefficient is the algebraic output
\begin{equation}\label{eq:c-from-B}
\boxed{c_n=-\frac{B_n}{1103872}.}
\end{equation}
Controller terminality therefore separates current-arity feedback from the autonomous memory subsystem.

\section{Triangular spectral extension and residual modules}\label{sec:spectral-extension}

We now prove that the factorization~\eqref{eq:low-poly} is an all-orders consequence of the fixed cubic contraction, rather than a finite-window coincidence.

Let $S$ denote the forward temporal shift,
\[
(Sx)(k)=x(k+1),
\]
and set
\[
E=S+2.
\]
Define the homogeneous operator
\[
A=-h\,\operatorname{ad}_\alpha.
\]
The $v^3$ recursion implies that
\[
g_3:=Ex_3
\]
is homogeneous:
\[
Sg_3=Ag_3.
\]
The exact cyclic space generated by $g_3(0)$ has dimension six.  In the basis
\[
g_3(0),Ag_3(0),\ldots,A^5g_3(0),
\]
the operator is
\begin{equation}\label{eq:Astar}
A_*=\begin{pmatrix}
0&0&0&0&0&384\\
1&0&0&0&0&32\\
0&1&0&0&0&-16\\
0&0&1&0&0&8\\
0&0&0&1&0&-4\\
0&0&0&0&1&2
\end{pmatrix}.
\end{equation}
Its characteristic and minimal polynomial are both
\begin{equation}\label{eq:mA}
\boxed{m_A(\lambda)=(\lambda+2)q(\lambda).}
\end{equation}

\begin{lemma}[Resonant Cauchy forcing]\label{lem:cauchy-resonance}
Let $a(k)=\lambda_0^ka_0$, let $B$ be a fixed bilinear map, and define
\[
C_b(k)=\sum_{i+j=k}B(a(i),b(j)).
\]
If $E=S-\lambda_0$, then
\[
(EC_b)(k)=B(a_0,b(k+1)).
\]
Consequently, $P(S)b=0$ implies $EP(S)C_b=0$.
\end{lemma}

\begin{proof}
Expand $C_b(k+1)$ and separate the term with $i=0$.  For $i\ge1$, use $a(i)=\lambda_0a(i-1)$ to identify the remaining sum with $\lambda_0C_b(k)$.  The final assertion follows because temporal polynomials commute with the fixed bilinear operation after one input is frozen to the geometric mode.
\end{proof}

Define the residual sequences
\[
r_b=m_A(S)x_b,
\qquad b=3,4,5,
\]
and their cyclic temporal modules
\[
\mathcal R_b=\C[S]r_b.
\]

\begin{theorem}[Residual Module Theorem]\label{thm:residual-modules}
For the fixed cubic contraction,
\[
\boxed{
\mathcal R_3\simeq \C[S]/(S+2),
\qquad
\mathcal R_4\simeq \C[S]/(S+2)^2,
}
\]
and
\[
\boxed{
\mu_{\mathcal R_5}(\lambda)=(\lambda+2)^3q(\lambda).
}
\]
Consequently
\[
\boxed{
p_5(\lambda)=(\lambda+2)^4q(\lambda)^2
}
\]
annihilates the $v^5$ rail for all temporal indices for which the autonomous triangular subsystem is defined.
\end{theorem}

\begin{proof}
Since $g_3=Ex_3$ is a cyclic homogeneous $A$-orbit with minimal polynomial $m_A=Eq$, one has
\[
r_3=q(S)g_3,
\qquad
Er_3=m_A(S)g_3=0.
\]
If $r_3=0$, the degree-five polynomial $q$ would annihilate the six-dimensional cyclic orbit $g_3$, contradicting~\eqref{eq:mA}.  Hence $\mathcal R_3\simeq\C[S]/E$.

Put $p_3=Em_A=E^2q$ and $y_4=p_3(S)x_4=Er_4$.  The term $[\xi,x_3]$ is killed by $p_3$, while Lemma~\ref{lem:cauchy-resonance} gives $E^2C_{22}=0$ for the $x_2*x_2$ convolution.  Thus $y_4$ is homogeneous.  Exact evaluation gives a nonzero eigenvector with
\[
Ay_4(0)=-2y_4(0),
\]
so $Ey_4=0$ but $y_4\neq0$.  Therefore $E^2r_4=0$ and $Er_4\neq0$, proving the second isomorphism.

Now put $p_4=Ep_3=E^3q$ and $y_5=p_4(S)x_5$.  The term $[\xi,x_4]$ is killed by $p_4$, and the resonant Cauchy lemma gives $p_4C_{23}=0$.  Hence $y_5$ is homogeneous.  Its first six orbit vectors are cyclic and carry exactly the same matrix~\eqref{eq:Astar}, so $\mu_{y_5}=m_A$.  It follows that $m_Ap_4$ annihilates $x_5$ and hence $p_4$ annihilates $r_5=m_Ax_5$.  Finally the eight cochains $r_5(0),\ldots,r_5(7)$ are linearly independent; an exact $8\times8$ rational minor has determinant
\[
-1974127506397904857128197160960\neq0.
\]
Since $\deg p_4=8$, no smaller polynomial annihilates $r_5$.  Thus $\mu_{\mathcal R_5}=p_4$.
\end{proof}

\begin{theorem}[Triangular Spectral Extension Lemma]\label{thm:triangular-spectral-extension}
Let $V$ be finite-dimensional, let $A\in\End(V)$, and let a temporal sequence satisfy
\[
Sx=Ax+f.
\]
Suppose $P$ is monic and $P(S)f=0$.  Set
\[
y=P(S)x
\]
and let
\[
d_y(\lambda)=\mu\bigl(A|_{\C[A]y(0)}\bigr).
\]
Then
\[
\mu_x\mid P\,d_y.
\]
If no polynomial of degree $<\deg P+\deg d_y$ annihilates $x$, then
\[
\boxed{\mu_x=P\,d_y.}
\]
In particular, if $y(0)$ is a nonzero $\lambda_0$-eigenvector, the new multiplier is $\lambda-\lambda_0$; if $y(0)$ is cyclic for the full homogeneous $A$-module, the new multiplier is $m_A$.
\end{theorem}

\begin{proof}
Since $P(S)$ commutes with the constant operator $A$,
\[
(S-A)y=P(S)(S-A)x=P(S)f=0.
\]
Thus $y(k)=A^ky(0)$, so $d_y(S)y=0$ and therefore $P(S)d_y(S)x=0$.  The equality statement follows from the stated lower-degree exclusion.
\end{proof}

The cubic transitions are therefore
\[
\boxed{
p_3=E\,m_A,
\qquad
p_4=E\,p_3,
\qquad
p_5=m_A\,p_4.
}
\]
They are, respectively, a full spectral extension, a simple resonant extension, and another full spectral extension.

\section{Residual landing operators and spectral selection}\label{sec:landing}

The preceding lemma can be made predictive using only finite boundary data.

\begin{proposition}[Finite-jet landing formula]\label{prop:finite-jet-landing}
Let
\[
x_{k+1}=Ax_k+f_k,
\qquad
P(z)=\sum_{r=0}^dp_rz^r,
\qquad
P(S)f=0.
\]
Then the homogeneous landing vector
\[
y_0=P(S)x(0)
\]
depends only on $x_0,f_0,\ldots,f_{d-1}$ and is
\[
\boxed{
y_0=P(A)x_0+\sum_{t=0}^{d-1}K_{P,t}(A)f_t,
}
\]
where
\[
K_{P,t}(z)=\sum_{r=t+1}^dp_rz^{r-1-t}.
\]
\end{proposition}

\begin{proof}
Unroll the recurrence to
\[
x_r=A^rx_0+\sum_{t=0}^{r-1}A^{r-1-t}f_t
\]
and substitute into $P(S)x(0)=\sum_rp_rx_r$.
\end{proof}

This defines a residual landing map $\Lambda_P$ from admissible finite lower-rail data to the homogeneous module.  The multiplier of a datum is the minimal polynomial of $A$ on the cyclic submodule generated by its landing vector.

\begin{theorem}[Complete landing operators for $x_4$ and $x_5$]\label{thm:complete-landing}
For the fixed cubic contraction, let $R=\C[S]$ and $E=S+2$.

For $x_4$,
\[
\mathscr D_4\simeq R/E^2R,
\qquad
\dim\mathscr D_4=2.
\]
In the basis $[1],[S]$ and the cyclic homogeneous basis of~\eqref{eq:Astar}, the complete landing operator is
\[
[\Lambda_4]=
\begin{pmatrix}
-192&384\\
80&-160\\
-32&64\\
12&-24\\
-4&8\\
1&-2
\end{pmatrix}.
\]
It satisfies
\[
\boxed{\rank\Lambda_4=1,
\qquad
(A+2I)\Lambda_4=0.}
\]
Hence every nonzero admissible $x_4$ landing is structurally an $E$-extension.

For $x_5$,
\[
\mathscr D_5\simeq R/p_4R,
\qquad
\dim\mathscr D_5=8,
\]
and, in the bases $[1],[S],\ldots,[S^7]$ and $y_5,Ay_5,\ldots,A^5y_5$,
\[
[\Lambda_5]=
\begin{pmatrix}
1&0&0&0&0&0&384&768\\
0&1&0&0&0&0&32&448\\
0&0&1&0&0&0&-16&0\\
0&0&0&1&0&0&8&0\\
0&0&0&0&1&0&-4&0\\
0&0&0&0&0&1&2&0
\end{pmatrix}.
\]
Thus
\[
\boxed{\rank\Lambda_5=6,}
\]
so $\Lambda_5$ is surjective onto the full homogeneous module.  Moreover
\[
\ker\Lambda_5=m_A\mathscr D_5
\]
is two-dimensional.
\end{theorem}

For $d=(d_0,\ldots,d_7)$ write
\[
r_d(S)=\sum_{j=0}^7d_jS^j.
\]
Then
\[
\Lambda_5(d)=r_d(A)y_5.
\]
The Krylov determinant satisfies
\begin{equation}\label{eq:landing-resultant}
\boxed{
\Delta_A(\Lambda_5(d))
=\det r_d(A)
=\operatorname{Res}(m_A,r_d).
}
\end{equation}
Since the datum $r_d=1$ gives determinant one, this polynomial is nonzero; full $m_A$-extension is therefore Zariski-generic.

\section{Spectral and extension stratifications}\label{sec:stratification}

The complete landing geometry is controlled by elementary gcd data in the PID $\C[S]$.

\begin{lemma}[GCD landing formula]\label{lem:gcd-landing}
Let $\mathcal H=\C[S]/(m)$ with cyclic generator $y=[1]$.  For $v=r(A)y$,
\[
\boxed{
\operatorname{Ann}_{\C[S]}(v)=\left(\frac{m}{\gcd(m,r)}\right).
}
\]
\end{lemma}

\begin{proof}
Write $m=gm_1$ and $r=gr_1$ with $g=\gcd(m,r)$ and $\gcd(m_1,r_1)=1$.  Then $fv=0$ iff $m\mid fr$, equivalently $m_1\mid fr_1$, and hence $m_1\mid f$.
\end{proof}

For the cubic homogeneous polynomial $m_A=Eq$, exact symbolic calculation gives
\[
q(-2)=-672\neq0,
\qquad
\gcd(q,q')=1,
\qquad
\gcd(m_A,m_A')=1,
\]
and $q$ is irreducible over $\Q$.  Over a splitting field write
\[
q(S)=\prod_{i=1}^5Q_i(S),
\qquad
Q_i(S)=S-\beta_i.
\]
Thus $m_A$ has six distinct spectral roots.

\begin{theorem}[Spectral landing stratification]\label{thm:spectral-stratification}
For $d\in\mathscr D_5$,
\[
\boxed{
\mu_H(d)=\frac{m_A}{\gcd(m_A,r_d)}.
}
\]
Over a splitting field, let
\[
Z(d)=\{\lambda\in\{-2,\beta_1,\ldots,\beta_5\}:r_d(\lambda)=0\}.
\]
Then the locally closed strata
\[
\mathscr S_Z=\{d:Z(d)=Z\}
\]
are all nonempty, there are $2^6=64$ of them, and on $\mathscr S_Z$
\[
\boxed{
\mu_H=\prod_{\lambda\notin Z}(S-\lambda).
}
\]
Moreover
\[
\dim\mathscr S_Z=8-|Z|
\]
and
\[
\overline{\mathscr S_Z}
=\bigsqcup_{Z'\supseteq Z}\mathscr S_{Z'}.
\]
\end{theorem}

Equation~\eqref{eq:landing-resultant} factors over $\Q$ as
\[
\operatorname{Res}(m_A,r_d)
=r_d(-2)\operatorname{Res}(q,r_d).
\]
Hence the rational exceptional divisor has a resonant linear component and an irreducible quintic norm component.  Rational data can therefore have only the four homogeneous multipliers
\[
m_A,
\qquad q,
\qquad E,
\qquad1.
\]

The forcing quotient retains higher resonant depth.  Set
\[
g_H=\gcd(m_A,r_d),
\qquad
g_F=\gcd(p_4,r_d).
\]

\begin{theorem}[Two-axis forcing/spectral stratification]\label{thm:two-axis}
Since $m_A\mid p_4$,
\[
\boxed{g_H=\gcd(m_A,g_F).}
\]
Over a splitting field define
\[
\nu(d)=\min\{3,\operatorname{ord}_{E}r_d\}
\]
and
\[
Z(d)=\{i:r_d(\beta_i)=0\}.
\]
Then the common refinement is indexed by
\[
(\nu,Z)\in\{0,1,2,3\}\times2^{\{1,\ldots,5\}},
\]
with
\[
\boxed{
g_F=E^\nu\prod_{i\in Z}Q_i,
\qquad
g_H=E^{\min(1,\nu)}\prod_{i\in Z}Q_i.
}
\]
There are $128$ nonempty strata, with
\[
\dim\mathscr S_{\nu,Z}=8-\nu-|Z|,
\]
and closure order
\[
(\nu,Z)\preceq(\nu',Z')
\iff
\nu\le\nu',\quad Z\subseteq Z'.
\]
The codimension generating polynomial is
\[
(1+t+t^2+t^3)(1+t)^5.
\]
\end{theorem}

The two exceptional resultants have the same reduced support but different resonant multiplicity:
\[
\operatorname{Res}(m_A,r_d)=r_d(-2)\operatorname{Res}(q,r_d),
\]
\[
\operatorname{Res}(p_4,r_d)=r_d(-2)^3\operatorname{Res}(q,r_d).
\]
Thus forcing memory is a resonant jet refinement of reduced spectral support.

Finally introduce the canonical higher-lift invariant
\[
g_X(d)=\gcd(p_5,r_d),
\qquad
p_5=E^4q^2.
\]
Here $r_d$ is always the canonical representative of degree $<8$ in $\mathscr D_5$.

\begin{theorem}[Three-level extension stratification]\label{thm:three-level}
For every canonical datum,
\[
\boxed{
g_H\mid g_F\mid g_X,
\qquad
g_H=\gcd(m_A,g_X),
\qquad
g_F=\gcd(p_4,g_X).
}
\]
Over a splitting field define
\[
a(d)=\min\{4,\operatorname{ord}_{E}r_d\}\in\{0,1,2,3,4\},
\]
and
\[
b_i(d)=\min\{2,\operatorname{ord}_{Q_i}r_d\}\in\{0,1,2\}.
\]
For a nonzero datum the exact contact label $(a,b_1,\ldots,b_5)$ is realizable if and only if
\[
\boxed{a+\sum_{i=1}^5b_i\le7.}
\]
On this stratum,
\[
\boxed{
g_X=E^a\prod_iQ_i^{b_i},
}
\]
while $g_F$ and $g_H$ are obtained by truncating exponents to $(3,1)$ and $(1,1)$ respectively.  The stratum has
\[
\boxed{
\dim=8-a-\sum_ib_i.
}
\]
There are $708$ nonzero fine strata and one zero stratum.  The canonical full-lift annihilator is
\[
\boxed{
\mu_X^{\mathrm{can}}=\frac{p_5}{g_X}
=E^{4-a}\prod_iQ_i^{2-b_i}.
}
\]
\end{theorem}

\begin{proof}
The divisibility chain follows from $m_A\mid p_4\mid p_5$.  Over the splitting field, gcds are determined by truncated contact orders at the six pairwise distinct spectral points.  A nonzero degree-$<8$ polynomial divisible by $E^a\prod_iQ_i^{b_i}$ exists exactly when the total contact weight is at most seven; sufficiency is realized by that product itself.  Hermite interpolation gives independence of the corresponding jet conditions and hence the dimension formula.  Summing feasible labels gives
\[
1+6+21+51+96+146+186+201=708.
\]
The zero datum forms the unique final stratum.
\end{proof}

Scheme-theoretically, the three levels have multiplicity profiles
\[
\operatorname{Res}(m_A,r_d):\quad (E,Q_i)=(1,1),
\]
\[
\operatorname{Res}(p_4,r_d):\quad (E,Q_i)=(3,1),
\]
\[
\operatorname{Res}(p_5,r_d):\quad (E,Q_i)=(4,2).
\]
The third level therefore adds precisely the fourth resonant jet and second-order contact with each quintic spectral hyperplane.

Over $\Q$, irreducibility of $q$ and the degree bound $\deg r_d<8$ imply that the nonzero possibilities for $g_X$ are exactly
\[
1,
E,
E^2,
E^3,
E^4,
q,
Eq,
E^2q.
\]
Thus the rational canonical three-level stratification consists of eight nonzero strata plus the zero datum.

\begin{remark}[Scope]\label{rem:strat-scope}
Theorem~\ref{thm:three-level} classifies the canonical polynomial lift in the fixed cubic $x_5$ module.  It is not a claim that every inverse-Leibniz contraction has the same multiplicity pattern.  The general problem is to determine the analogue of the exponent truncation rules from the spectral multiplicity vectors of the homogeneous operator and the triangular forcing maps.
\end{remark}

\section{The central-commutator plane as a future application}

The motivating algebraic example is the central-commutator algebra
\[
A_{\mathrm{cc}}
=
\C\langle z,\bar z,J\rangle
\Big/
\bigl([z,\bar z]-2J,[J,z],[J,\bar z]\bigr),
\]
optionally equipped with
\[
z^*=\bar z,
\qquad
J^*=J.
\]
The basic identity
\[
[z,\bar z]=-2i[x,y],
\qquad
z=x+iy,
\]
then converts the relation $[z,\bar z]=2J$ into a central real-coordinate commutator.  The present paper does not construct a geometric realization of this algebra.  Instead, it supplies the reconstruction and obstruction machinery needed before one can consistently ask whether a bilinear geometric datum $B(z,\bar z)$ determines compatible algebra actions.

\section{Open problems and next steps}

\begin{problem}[Relative presentation versus intrinsic normal forms]
Theorem~\ref{thm:strict-transfer-no-go} gives a nonzero presentation quartic class, while Theorem~\ref{thm:intrinsic-ell4-vanish} shows that it disappears after restoring the full Hochschild differential.  Construct a relative theory, preferably in terms of the mapping cone or homotopy fiber of
\[
\Psi:\ghoch_{B,p}\to\gpres_{B,p},
\]
that measures exactly which presentation normal-form obstructions are truncation artifacts and which survive intrinsically.
\end{problem}

\begin{problem}[General triangular spectral-extension theory]
The fixed cubic contraction now admits an all-orders low-rail theorem and complete landing stratifications.  Abstract these results to triangular systems whose homogeneous operator has spectral multiplicity data
\[
m_A(S)=\prod_\lambda(S-\lambda)^{e_\lambda}
\]
and whose forcing branches are built from linear maps and Cauchy convolutions of lower rails.  Determine general conditions under which forcing annihilators, residual landing modules, and higher-lift contact are obtained by deterministic exponent-truncation rules analogous to Theorems~\ref{thm:triangular-spectral-extension} and~\ref{thm:three-level}.
\end{problem}

\begin{problem}[Stationary reuse and reserve actuators]
Prove or disprove that the finite box $U_{16}$ remains sufficient at all subsequent full harmonic arities beyond the verified stationary regime.  If the stationary phase terminates, determine whether the reserved two-sided direction $E_{33}$ is the next actuator or whether a new support-level bifurcation occurs.  Relate this question to the all-orders behavior of the low-rail forcing observable~\eqref{eq:B-low-observable}.
\end{problem}

\begin{problem}[Minimality of a cubic obstruction]
The example of Section~\ref{sec:cubic-example} uses $A=\C[\epsilon]/(\epsilon^4)$ and $\dim V=3$.  Determine whether a genuinely cubic inverse-Leibniz obstruction can occur for a smaller algebra or target dimension.  The bounded truncated-polynomial Jordan search used in this draft found none for the smaller cases tested, but this is not a proof of global minimality.
\end{problem}

\begin{problem}[Effective relation kernels and presentation syzygies]
Study the map
\[
\overline\Pi_{\mathrm{eff}}:\Obs^{\mathrm{eff}}_p\longrightarrow\coker\cD_p
\]
in larger examples.  Distinguish intrinsic effective classes killed by this map from excess directions in $\coker\cD_p$ caused by redundant equations or higher syzygies.
\end{problem}

\begin{problem}[Gauge quotient]
Restore degree zero and study reconstruction up to simultaneous change of basis in $V$.  Identify the corresponding Deligne groupoid or derived moduli problem, and compare framed rigidity with rigidity modulo conjugation.
\end{problem}

\begin{problem}[Global symmetry and identifiability]
Characterize finite groups or groupoids acting on $\Rfib(B)$ without changing $B$.  The four-point example suggests coordinatewise left-right swaps as a discrete ambiguity invisible to local tangent data.
\end{problem}

\begin{problem}[Graded reconstruction]
For finitely generated graded algebras, construct finite-level observability matrices and identify compatibility conditions ensuring that solutions extend through all word lengths.
\end{problem}

\begin{problem}[Quotient obstructions]
Determine how relations in $A=\C\langle z_1,\ldots,z_s\rangle/I$ constrain weak actions and create new obstruction classes beyond those visible in the free algebra.
\end{problem}

\begin{problem}[Literature positioning]
Compare the inverse Leibniz problem systematically with universal differential calculi, Hochschild cohomology, double derivations, representation schemes, and deformation theory of modules and algebra homomorphisms.  No novelty claim should be finalized before this comparison is complete.
\end{problem}

\appendix

\section{Exact computational verification}

The accompanying scripts perform the calculations over exact integers and rationals.  The early scripts verify the reconstruction examples, relation morphism, effective obstruction quotients, cubic Kuranishi obstruction, complete formal germ, presentation quartic no-go theorem, and intrinsic Hochschild transfer.  The later verifier family continues the same fixed intrinsic calculation through the finite-feedback, block-actuator, stationary-reuse, low-rail, residual-module, landing-operator, and stratification stages.  Full harmonic feedback is verified through arity twenty-two, while the autonomous low-rail recursion is independently extended through arity twenty-three.  Modular sparse minors are used only to certify rational lower bounds when an independent exact rational argument supplies the matching upper bound; all theorem statements in the later sections record this distinction explicitly.  For the four-point example it verifies:
\[
M_B\in M_{256\times32}(\mathbb Q),
\qquad
\rank M_B=28,
\qquad
\dim\Wfib(B)=4.
\]
After substituting the weak parameterization, the strict Groebner basis is
\[
(s_{11}-1)(s_{11}-27),\quad s_{12},\quad s_{21},\quad(s_{22}-8)(s_{22}-64).
\]
The Jacobian of the complete strict system has rank $32$ at each of the four points, so every strict tangent space is zero.

For the dual-number example, the exact Groebner basis is
\[
\ell+r,\qquad r^2,
\]
and the Jacobian at the origin has rank $1$, yielding a one-dimensional tangent space.  Direct expansion over $\C[t]/(t^3)$ verifies that every nonzero tangent vector produces nonzero second-order squares.

The same script constructs the normalized constrained Hochschild complex for $S=A^e$ at the origin.  It verifies
\[
\dim\ghoch^1_{B,p}=2,
\qquad
\dim Z^1(\ghoch_{B,p})=1,
\]
\[
\rank(d_2)=5,
\qquad
\dim Z^2(\ghoch_{B,p})=4,
\qquad
\dim H^2(\ghoch_{B,p})=3.
\]
It also checks that the cup square of the tangent generator is closed but not in the image of $d_1$ restricted to $\ghoch^1_{B,p}$, and verifies the differential graded Leibniz identity on all normalized basis one-cochains.  The v0.3 script additionally verifies the chain-map and bracket identities of Proposition~\ref{prop:relation-dgla-morphism}, constructs the explicit basis
\[
H^2=\C[\eta]\oplus\C[\eta_1]\oplus\C[\eta_2],
\]
computes
\[
\rank\bigl(H^2(\ghoch_{B,p})\to\coker\cD_p\bigr)=3,
\]
and confirms the transferred Maurer--Cartan equation $t^2\eta=0$.  The v0.4 script further verifies
\[
\dim\Obs^{\mathrm{eff}}_p=1
\]
for the dual-number germ.  For the square-zero two-parameter example it verifies
\[
\dim\Wfib(B)_{\mathrm{unit}}=2,
\qquad
\dim\ghoch^1_{B,p}=18,
\qquad
\dim H^1=2,
\]
\[
\dim Z^2=64,
\qquad
\dim H^2=48,
\qquad
\dim\Obs^{\mathrm{eff}}_p=2,
\]
together with the exact strict ideal
\[
(uv,v^2),
\]
the independence of $[\eta_{uv}]$ and $[\eta_{vv}]$ modulo boundaries, and the rank-two effective relation map.

For the cubic example of Section~\ref{sec:cubic-example}, the v0.5 script verifies
\[
\dim\mathcal K_B=9,
\qquad
\rank(\cD_p|_{\mathcal K_B})=7,
\qquad
\dim T_p\mathbf R_B=2.
\]
It checks the displayed tangent basis, the weak second-order corrections $\eta$ and $\theta$, the vanishing of every quadratic Kuranishi product on $H^1$, and the exact identities
\[
\cD_p\eta=-\cQ_p(\xi),
\qquad
\cD_p\theta=-\cP_p(\alpha,\eta).
\]
Finally, it constructs the sparse functional \eqref{eq:cubic-certificate} and verifies
\[
\Lambda\circ\cD_p|_{\mathcal K_B}=0,
\qquad
\Lambda\bigl(\cP_p(\xi,\eta)\bigr)=8,
\]
which is an exact certificate for the nonzero cubic obstruction.


The v0.6 extension substitutes the full nine-parameter weak fiber into all strict equations and computes, over $\Q$, the exact lexicographic Groebner basis of Theorem~\ref{thm:cubic-complete-germ}.  It then verifies the rational formal section
\[
e=-\frac{3v^2}{1+2u}
\]
and checks directly that every remaining strict residual is divisible by $v^3$, with at least one unit quotient.  Finally it verifies
\[
\frac18\Lambda(\text{residual})
=
\frac{v^3\bigl(v+4(1+2u)\bigr)}{4(1+2u)^2}
\]
and the exact unit normalization to $v^3$.

The v0.7 extension verifies the sparse quartic certificate $\Gamma$, the rank-three image of $\mu_\alpha$, and the inconsistent presentation contraction system with coefficient rank $20$ and augmented rank $21$.  The v0.8 extension constructs the normalized constrained Hochschild degree-one space for the cubic example and verifies
\[
\dim\ghoch^1_{B,p}=90,
\qquad
\rank d_1=88.
\]
It checks the exact cochain identities
\[
\mathsf a\smile\mathsf a=0,
\qquad
[\mathsf a,\mathsf x]=0,
\qquad
 d_\mu\mathsf e=-\mathsf x\smile\mathsf x,
\qquad
[\mathsf a,\mathsf e]=-2\mathsf x\smile\mathsf x,
\]
constructs the explicit cochain $\mathsf k$, and verifies that
\[
\widehat\zeta=[\mathsf x,\mathsf e]+d_\mu\mathsf k
\]
is closed and maps to the nonzero cubic presentation class.  Finally, for the quartic sources $Y_{13},Y_{04}$ it verifies the exact degree-three minor
\[
\begin{pmatrix}
0&-\tfrac12\\
\tfrac12&\tfrac12
\end{pmatrix}
\]
with determinant $1/4$, together with
\[
\Gamma(\Psi^2Y_{13})=-16,
\qquad
\Gamma(\Psi^2Y_{04})=-\frac72.
\]
These are the exact certificates used in Theorem~\ref{thm:intrinsic-ell4-vanish} and Proposition~\ref{prop:intrinsic-presentation-gap}.

The v0.9 extension continues the same intrinsic contraction through arity five.  It verifies
\[
\rank\Span\{d\mathsf e,d\mathsf k,d\mathsf p,d\mathsf q\}=4
\]
and constructs the quintic sources \eqref{eq:arity5-Z32}--\eqref{eq:arity5-Z05}.  For the joint nonclosed source space it verifies the exact degree-three minor
\[
\begin{pmatrix}
0&-\tfrac12&0&-\tfrac23&-\tfrac12\\
\tfrac12&\tfrac12&-\tfrac13&\tfrac13&-1\\
0&1&0&\tfrac13&0\\
0&0&0&\tfrac13&0\\
\tfrac12&-\tfrac12&\tfrac13&-\tfrac13&0
\end{pmatrix},
\qquad
\det=\frac1{18}.
\]
Thus the same contraction can be chosen with both $\ell_4$ and $\ell_5$ vanishing on the $H^1$-generated sector.

The v0.10 extension continues that contraction through arity six.  It verifies
\[
\rank\Span\{d\mathsf e,d\mathsf k,d\mathsf p,d\mathsf q,d\mathsf r,d\mathsf w,d\mathsf t\}=7,
\]
checks the five sextic source formulas \eqref{eq:arity6-U42}--\eqref{eq:arity6-U06}, and verifies the boundary identity \eqref{eq:arity6-U42-boundary}.  On the joint arity-four-through-six nonclosed source space, the exact degree-three certificate \eqref{eq:arity6-minor} has determinant $-1/486$.  Hence the joint differential rank is nine and one intrinsic contraction can be chosen with $\ell_4$, $\ell_5$, and $\ell_6$ all vanishing on the $H^1$-generated sector.

The v0.11 extension continues the same contraction through arity seven.  It verifies the four pure-mixed sextic homotopy lifts $\mathsf A,\mathsf B,\mathsf C,\mathsf D$, the rank-eleven degree-one splitting \eqref{eq:arity7-degree1-rank}, all six arity-seven source formulas \eqref{eq:arity7-V52}--\eqref{eq:arity7-V07}, and the boundary identity \eqref{eq:arity7-V52-boundary}.  On the joint arity-four-through-seven nonclosed source space, the exact degree-three certificate \eqref{eq:arity7-minor} has determinant $11/34992$.  Hence the joint differential rank is fourteen and one intrinsic contraction can be chosen with $\ell_4$, $\ell_5$, $\ell_6$, and $\ell_7$ all vanishing on the $H^1$-generated sector.



For the higher intrinsic calculation, the exact verifier records the following checkpoints.  The finite-box failures at arities thirteen through eighteen occur at the $\eta$-coordinates displayed in Table~\ref{tab:finite-box-bifurcation}; the fresh mixed directions in that table restore acyclic transfer.  The block-actuator calculation proves the kernel flag of Theorem~\ref{thm:block-actuator-flag} and the exact relation $T_\alpha(E_{33})=T_\xi(E_{23})$.  The rescued source-history ranks through arities nineteen, twenty, twenty-one, and twenty-two are respectively
\[
126,\qquad141,\qquad157,\qquad174.
\]
The fixed box $U_{16}$ gives full-row-rank harmonic systems of sizes $28\times455$, $30\times490$, $32\times525$, and $34\times560$ at those arities.  The low-rail verifier checks every available transition of equations~\eqref{eq:x2-rec}--\eqref{eq:x5-rec}, the frozen recurrence~\eqref{eq:Y14-rec}, and two independent out-of-sample predictions at $Y_{17}$ and $Y_{18}$.  The residual-module verifier constructs the cyclic matrix~\eqref{eq:Astar}, verifies $m_A=Eq$, proves the exact Jordan residual modules of Theorem~\ref{thm:residual-modules}, and records the nonzero $8\times8$ minor used for $\mu_{\mathcal R_5}=p_4$.  The landing-operator verifier checks both matrices in Theorem~\ref{thm:complete-landing}, including $(A+2I)\Lambda_4=0$, surjectivity of $\Lambda_5$, and the resultant identity~\eqref{eq:landing-resultant}.  Finally, the stratification verifiers check squarefreeness and irreducibility of $q$, all gcd truncation identities, the $64$-, $128$-, and $709$-stratum counts, the Hermite-jet rank certificates, and the rational reductions described in Section~\ref{sec:stratification}.

\begin{thebibliography}{99}

\bibitem{Gerstenhaber1964}
M.~Gerstenhaber,
\emph{On the deformation of rings and algebras},
Ann. of Math. (2) \textbf{79} (1964), 59--103.
\href{https://doi.org/10.2307/1970484}{doi:10.2307/1970484}.

\bibitem{Schlessinger1968}
M.~Schlessinger,
\emph{Functors of Artin rings},
Trans. Amer. Math. Soc. \textbf{130} (1968), 208--222.
\href{https://doi.org/10.1090/S0002-9947-1968-0217093-3}{doi:10.1090/S0002-9947-1968-0217093-3}.

\bibitem{CuntzQuillen1995}
J.~Cuntz and D.~Quillen,
\emph{Algebra extensions and nonsingularity},
J. Amer. Math. Soc. \textbf{8} (1995), 251--289.
\href{https://doi.org/10.1090/S0894-0347-1995-1303029-0}{doi:10.1090/S0894-0347-1995-1303029-0}.

\bibitem{VanDenBergh2008}
M.~Van den Bergh,
\emph{Double Poisson algebras},
Trans. Amer. Math. Soc. \textbf{360} (2008), 5711--5769.
\href{https://doi.org/10.1090/S0002-9947-08-04518-2}{doi:10.1090/S0002-9947-08-04518-2}.


\bibitem{StacksCotangent}
The Stacks Project Authors,
\emph{The cotangent complex and the naive cotangent complex},
Tags \href{https://stacks.math.columbia.edu/tag/00S0}{00S0} and
\href{https://stacks.math.columbia.edu/tag/08R6}{08R6}.

\end{thebibliography}

\end{document}
